% Complete reviewed Unit 11; editorial master: chapters/ch11/student.md
% Markdown SHA-256: b134bcfe29fae274bf244cd22463d1b24bcc428814066dfcac174a2a673ab527
% One-use initial assembly followed by complete content and page review.
\newcommand{\runningtitle}{Unit 11. Adverbial clauses}
% Layout inherited from the reviewed previous course; set a running title per unit.
% This is a small hand-maintained TeX layout, not a conversion pipeline.
\RequirePackage{fix-cm}
\documentclass[11pt,twoside]{article}
\usepackage{fontspec}
\setmainfont{Cambria}[
  Renderer=HarfBuzz,
  BoldFont=cambriab.ttf,
  ItalicFont=cambriai.ttf,
  BoldItalicFont=cambriaz.ttf
]
\setmonofont{lmmono10-regular.otf}[BoldFont=lmmonolt10-bold.otf]
\usepackage[paperwidth=165mm,paperheight=241mm,inner=21mm,outer=19mm,
  top=22mm,bottom=20mm,headheight=14pt,headsep=7mm,footskip=10mm]{geometry}
\usepackage{xcolor,array,booktabs,longtable,tabularx}
\usepackage{enumitem,fancyhdr}
\usepackage[unicode,hidelinks]{hyperref}
\definecolor{ink}{HTML}{000000}
\definecolor{muted}{HTML}{4C5558}
\setlength{\parindent}{0pt}
\setlength{\parskip}{3.5pt}
\AtBeginDocument{\fontsize{11}{14.3}\selectfont}
\setlength{\tabcolsep}{4pt}
\renewcommand{\arraystretch}{1.10}
\setlist[itemize]{leftmargin=1.3em,label=\textbullet,itemsep=2pt,topsep=3pt,parsep=0pt}
\setlist[enumerate]{leftmargin=1.8em,itemsep=3pt,topsep=4pt,parsep=0pt}
\widowpenalty=10000
\clubpenalty=10000
\emergencystretch=1em
\raggedbottom
\pagestyle{fancy}
\fancyhf{}
\providecommand{\runningtitle}{Old English}
\fancyhead[LE,RO]{\small\itshape \runningtitle}
\fancyhead[LO,RE]{\small Old English}
\fancyfoot[LE,RO]{\small\thepage}
\renewcommand{\headrulewidth}{0.3pt}
\newcommand{\booktitle}[2]{
  {\fontsize{9}{11}\selectfont\color{muted}#1\par}
  \vspace{3pt}
  {\fontsize{20}{24}\selectfont\color{ink}#2\par}
  \vspace{6pt}
}
\newcommand{\source}[1]{\par{\footnotesize\color{muted}#1\par}}
\newcommand{\readingtext}[1]{
  \par\addvspace{5pt}
  \begingroup\leftskip=1.5em
    \fontsize{12}{16}\selectfont #1\par
  \endgroup\addvspace{5pt}
}
\makeatletter
\renewcommand{\section}{\@startsection{section}{1}{\z@}
  {-2.4ex \@plus -0.6ex \@minus -0.2ex}{0.8ex \@plus 0.2ex}
  {\normalfont\fontsize{12}{15}\selectfont\bfseries\color{ink}}}
\renewcommand{\subsection}{\@startsection{subsection}{2}{\z@}
  {-1.8ex \@plus -0.5ex \@minus -0.2ex}{0.5ex \@plus 0.2ex}
  {\normalfont\fontsize{11}{14}\selectfont\bfseries\color{ink}}}
\makeatother

\hyphenpenalty=10000
\exhyphenpenalty=10000
\emergencystretch=3em
\newcommand{\itempara}[2]{\par\addvspace{3pt}\noindent\hangindent=1.8em\hangafter=1\makebox[1.8em][l]{#1.}#2\par}
\newcommand{\reserve}[1]{\par\penalty-100\begingroup\dimen0=\pagegoal\advance\dimen0 by-\pagetotal\ifdim\dimen0<#1\newpage\fi\endgroup}
\begin{document}
\addfontfeatures{Ligatures=TeXOff}
\raggedright
\booktitle{UNIT 11}{Adverbial clauses}

% mdparagraph:P001:begin
An adverbial clause places an action in time or relates it to a cause, purpose, result, concession or condition. First identify that relation; then consider how the speaker presents it. An indicative commonly asserts an event. A subjunctive can present an intention, an anticipated event, a concession or a supposition. It does not always mean that something is doubtful or false.
% mdparagraph:P001:end

% mdparagraph:P002:begin
The fourteen translated examples, the exercises marked F, A, T, R, B and S, and “The fisherman and his son” are new teaching material. Cook IV.16–17 and V.18 retain their original English prompts. The two later readings are exact passages from Cook’s \emph{First Book}, with separate context, help and original notes.
% mdparagraph:P002:end

% mdparagraph:P003:begin
In the tables, 1, 2 and 3 denote person. Repeated spellings still occupy distinct grammatical cells. Nominative, genitive, dative and accusative forms are for use in composition within their stated meanings. Forms marked \textbf{reference} are for recognition. Supplied expressions may be used with their printed help. Preserve vowel lengths and the distinction between e and ę.
% mdparagraph:P003:end

\reserve{7\baselineskip}
% mdparagraph:P004:begin
\section*{1. Time clauses}
% mdparagraph:P004:end

% teaching: NEW-U11-GRAM-u11-clause-roles-order
% teaching: NEW-U11-GRAM-u11-past-when-then
% teaching: NEW-U11-VOC-THA-CONJUNCTION
% teaching: NEW-U11-VOC-THA-ADVERB
% mdparagraph:P005:begin
For an event at a particular past time, \textbf{ðā}, “when”, introduces a subordinate clause, normally with the indicative. A second \textbf{ðā}, “then”, may introduce the main clause. The first connects clauses; the second is an adverb inside the main clause. The same spelling also belongs to familiar inflected forms of the article or demonstrative: identify its function from the construction.
% mdparagraph:P005:end

\reserve{8\baselineskip}
% example: NEW-U11-EXAMPLE01
% mdparagraph:P006:begin
Ðā se fiscere cōm, ðā sæt se sunu on ðǣm hūse.
% mdparagraph:P006:end

% mdparagraph:P007:begin
When the fisherman came, the son was sitting in the house.
% mdparagraph:P007:end

% mdparagraph:P008:begin
Cōm and sæt are past indicative forms. The first clause tells when the sitting took place. The second ðā resumes the narrative; English can omit “then”.
% mdparagraph:P008:end

% mdparagraph:P009:begin
A dependent clause often puts its finite verb late. A main clause beginning with an adverb often puts its finite verb before its subject, as in ðā sæt se sunu. These are useful tendencies, not instructions to reverse or move every verb mechanically. Case and agreement still determine the participants.
% mdparagraph:P009:end

% teaching: NEW-U11-GRAM-u11-recurring-when-then
% teaching: NEW-U11-VOC-THONNE-CONJUNCTION
% teaching: NEW-U11-VOC-THONNE
% mdparagraph:P010:begin
Use \textbf{ðonne}, “when”, for recurring or prospective time. It too can be answered by the adverb \textbf{ðonne}, “then”. In the following example the visit and gift recur. Both finite verbs are present indicative. A future setting does not by itself require a subjunctive or a separate future ending.
% mdparagraph:P010:end

\reserve{8\baselineskip}
% example: NEW-U11-EXAMPLE02
% mdparagraph:P011:begin
Ðonne se munuc cymð, ðonne sęlð se mann him męte.
% mdparagraph:P011:end

% mdparagraph:P012:begin
When the monk comes, the man gives him food.
% mdparagraph:P012:end

% teaching: NEW-U11-GRAM-u11-while
% teaching: NEW-U11-VOC-THA-HWILE-THE
% mdparagraph:P013:begin
\textbf{Ðā-hwīle-ðe}, “while”, introduces a period during which another action takes place. In a statement of actual simultaneous events, use the indicative.
% mdparagraph:P013:end

\reserve{8\baselineskip}
% example: NEW-U11-EXAMPLE03
% mdparagraph:P014:begin
Ðā-hwīle-ðe se munuc sang, se sunu dranc wæter.
% mdparagraph:P014:end

% mdparagraph:P015:begin
While the monk sang, the son drank water.
% mdparagraph:P015:end

% mdparagraph:P016:begin
The two actions overlap in time. The conjunction alone does not make either action the cause of the other.
% mdparagraph:P016:end

% teaching: NEW-U11-GRAM-u11-mid-thy-temporal
% teaching: NEW-U11-VOC-MID-THY
% mdparagraph:P017:begin
\textbf{Mid-ðȳ}, “when; while”, is another temporal conjunction. Learn the whole connection in this use; do not translate it word by word as a phrase of means.
% mdparagraph:P017:end

\reserve{8\baselineskip}
% example: NEW-U11-EXAMPLE04
% mdparagraph:P018:begin
Mid-ðȳ se sunu ǣt, se fæder sæt on ðǣm hūse.
% mdparagraph:P018:end

% mdparagraph:P019:begin
While the son ate, the father sat in the house.
% mdparagraph:P019:end

% teaching: NEW-U11-GRAM-u11-before-anticipated
% teaching: NEW-U11-VOC-AER-CONJUNCTION
% mdparagraph:P020:begin
\textbf{Ǣr}, “before”, can introduce an event anticipated from the viewpoint of the main clause. After the following instruction, the monk’s coming is still expected, so \textbf{cume} is present subjunctive. The command itself is plural.
% mdparagraph:P020:end

\reserve{8\baselineskip}
% example: NEW-U11-EXAMPLE05
% mdparagraph:P021:begin
Ne farað, ǣr se munuc cume.
% mdparagraph:P021:end

% mdparagraph:P022:begin
Do not go before the monk comes.
% mdparagraph:P022:end

% teaching: NEW-U11-GRAM-u11-until-anticipated
% teaching: NEW-U11-VOC-OTH-THAET
% teaching: NEW-U11-VOC-HER
% mdparagraph:P023:begin
\textbf{Oð-ðæt}, “until”, gives the endpoint of an action or state. Use the subjunctive for the selected anticipated endpoint. \textbf{Hēr} means “here”.
% mdparagraph:P023:end

\reserve{8\baselineskip}
% example: NEW-U11-EXAMPLE06
% mdparagraph:P024:begin
Wes hēr, oð-ðæt wē cumen.
% mdparagraph:P024:end

% mdparagraph:P025:begin
Be here until we come.
% mdparagraph:P025:end

% mdparagraph:P026:begin
Wes addresses one person. Cumen is present subjunctive plural: the people expected to come are wē. The English present “come” does not decide the Old English mood.
% mdparagraph:P026:end

% mdparagraph:P027:begin
These are statements about viewpoint, not universal rules equating “before” or “until” with one tense or mood. A narrator can use a subjunctive for an event viewed as still ahead at an earlier point, even if it later happened.
% mdparagraph:P027:end

\reserve{7\baselineskip}
% mdparagraph:P028:begin
\section*{2. Cause, purpose and result}
% mdparagraph:P028:end

% teaching: NEW-U11-GRAM-u11-causal-for-thon-the
% teaching: NEW-U11-VOC-FOR-THON-THE
% mdparagraph:P029:begin
\textbf{For-ðon-ðe}, “because”, introduces a stated reason. Use the indicative for the factual cause illustrated here. The earlier causal for-ðām remains available; the exact reading below also has shorter for-ðon.
% mdparagraph:P029:end

\reserve{8\baselineskip}
% example: NEW-U11-EXAMPLE07
% mdparagraph:P030:begin
Se mann healp ðǣm cuman, for-ðon-ðe hē wæs untrum.
% mdparagraph:P030:end

% mdparagraph:P031:begin
The man helped the stranger because the stranger was sick.
% mdparagraph:P031:end

% mdparagraph:P032:begin
Here hē is understood to mean the stranger. Masculine singular grammar by itself could also fit the man; the intended reference must be established from context. The reason clause explains why the help was given.
% mdparagraph:P032:end

% teaching: NEW-U11-GRAM-u11-purpose
% teaching: NEW-U11-VOC-THAET-CONJUNCTION
% mdparagraph:P033:begin
\textbf{Ðæt + subjunctive} introduces an intended purpose: “so that; in order that”. Unlike the earlier same-subject purpose infinitive, a finite purpose clause has its own expressed or understood subject. That subject may be the person acting in the main clause or a different person.
% mdparagraph:P033:end

\reserve{8\baselineskip}
% example: NEW-U11-EXAMPLE08
% mdparagraph:P034:begin
Ic cume, ðæt ic ðē helpe.
% mdparagraph:P034:end

% mdparagraph:P035:begin
I come so that I may help you.
% mdparagraph:P035:end

% mdparagraph:P036:begin
The first cume is indicative in the statement “I come”; helpe is subjunctive in the purpose. First-person singular present forms may have the same spelling in both moods. Function and context matter as well as endings.
% mdparagraph:P036:end

\reserve{8\baselineskip}
% example: NEW-U11-EXAMPLE09
% mdparagraph:P037:begin
Se fæder cōm, ðæt hē his suna hulpe.
% mdparagraph:P037:end

% mdparagraph:P038:begin
The father came so that he might help his son.
% mdparagraph:P038:end

% mdparagraph:P039:begin
Hulpe is third singular preterite subjunctive of helpan. The father’s purpose is to help his son; the sentence does not by itself assert that he succeeded. His suna is dative, the person helped, with his as genitive possessor.
% mdparagraph:P039:end

% mdparagraph:P040:begin
The present purpose example uses helpe; the past example uses hulpe. For the already learned subjunctives, present singular -e contrasts with plural -en; the corresponding strong preterite uses the preterite plural stem with singular -e and plural -en. Thus helpe/\allowbreak{}helpen and hulpe/\allowbreak{}hulpen differ in tense and number. The examples here establish these particular present and past settings, without a rule mechanically replacing every subordinate tense.
% mdparagraph:P040:end

% teaching: NEW-U11-GRAM-u11-negative-purpose
% teaching: NEW-U11-VOC-THY-LAES-THE
% mdparagraph:P041:begin
\textbf{Ðȳ-lǣs-ðe + subjunctive}, “lest; so that … not”, states an outcome to be prevented. Its negative force belongs to the purpose clause. It does not negate the main instruction.
% mdparagraph:P041:end

\reserve{8\baselineskip}
% example: NEW-U11-EXAMPLE10
% mdparagraph:P042:begin
Heald ðæt dēor, ðȳ-lǣs-ðe hit ðone męte nime.
% mdparagraph:P042:end

% mdparagraph:P043:begin
Hold the animal, lest it take the food.
% mdparagraph:P043:end

% mdparagraph:P044:begin
Nime is third singular present subjunctive. The purpose is that the animal should not take the food. No actual taking is asserted.
% mdparagraph:P044:end

% teaching: NEW-U11-GRAM-u11-factual-result
% mdparagraph:P045:begin
An asserted result commonly has \textbf{swā … ðæt + indicative}, “so … that”. Swā marks the degree; the ðæt clause says what followed. Compare intended help with an actual consequence.
% mdparagraph:P045:end

\reserve{8\baselineskip}
% example: NEW-U11-EXAMPLE11
% mdparagraph:P046:begin
Se sunu wæs swā glæd, ðæt hē sang.
% mdparagraph:P046:end

% mdparagraph:P047:begin
The son was so glad that he sang.
% mdparagraph:P047:end

% mdparagraph:P048:begin
Sang is past indicative: the son did sing. A purpose clause would present singing as an aim. The earlier content clause, as after witan “know”, instead states what is known. The conjunction ðæt alone does not distinguish these three functions.
% mdparagraph:P048:end

% mdparagraph:P049:begin
Cook also gives a result-like optative curse under his subjunctive examples. That special wish does not make every factual result subjunctive. Read the force of the whole sentence.
% mdparagraph:P049:end

\reserve{7\baselineskip}
% mdparagraph:P050:begin
\section*{3. Concession, condition and comparison}
% mdparagraph:P050:end

% teaching: NEW-U11-GRAM-u11-concession
% teaching: NEW-U11-VOC-THEAH
% mdparagraph:P051:begin
\textbf{Ðēah}, “although; even though”, introduces a concession: a circumstance that might have been expected to prevent the main action. The selected construction normally has the subjunctive. The speaker can nevertheless accept the circumstance as a fact.
% mdparagraph:P051:end

\reserve{8\baselineskip}
% example: NEW-U11-EXAMPLE12
% mdparagraph:P052:begin
Ðēah se fiscere sīe eald, hē hilpð his suna.
% mdparagraph:P052:end

% mdparagraph:P053:begin
Although the fisherman is old, he helps his son.
% mdparagraph:P053:end

% mdparagraph:P054:begin
Sīe is third singular present subjunctive. The example concedes the fisherman’s old age; it does not deny it. Hilpð is indicative, asserting the help.
% mdparagraph:P054:end

% teaching: NEW-U11-GRAM-u11-contrary-present-condition
% mdparagraph:P055:begin
The earlier open \textbf{gif}, “if”, construction leaves the condition open. An explicitly \textbf{contrary-present} condition imagines something known not to be so now. In that setting both clauses use preterite subjunctive forms. Their past morphology need not give the sentence past-time reference.
% mdparagraph:P055:end

\reserve{12\baselineskip}
% mdparagraph:P056:begin
Open condition: the speaker does not know whether the monk is here.
% mdparagraph:P056:end

\reserve{8\baselineskip}
% example: NEW-U11-EXAMPLE13
% mdparagraph:P057:begin
Gif se munuc is hēr, hē hilpð ðǣm męnn.
% mdparagraph:P057:end

% mdparagraph:P058:begin
If the monk is here, he helps the man.
% mdparagraph:P058:end

\reserve{12\baselineskip}
% mdparagraph:P059:begin
Contrary-present condition: the monk is absent now.
% mdparagraph:P059:end

\reserve{8\baselineskip}
% example: NEW-U11-EXAMPLE14
% mdparagraph:P060:begin
Gif se munuc wǣre hēr, hē hulpe ðǣm męnn.
% mdparagraph:P060:end

% mdparagraph:P061:begin
If the monk were here, he would help the man.
% mdparagraph:P061:end

% mdparagraph:P062:begin
Wǣre and hulpe are preterite subjunctives referring to the imagined present. Do not translate them as completed past events merely because of their forms. This contrast does not introduce the full system of contrary-to-fact past conditions.
% mdparagraph:P062:end

% teaching: NEW-U11-GRAM-u11-ac-coordination
% teaching: NEW-U11-VOC-AC
% mdparagraph:P063:begin
\textbf{Ac}, “but”, coordinates contrasting clauses or predicates. It imposes no subordinate-clause mood. Keep the persons and numbers of both subjects even when English repeats “shall” or “will”; the already learned present can refer to the future.
% mdparagraph:P063:end

% teaching: NEW-U11-GRAM-u11-swa-comparative-clause
% teaching: NEW-U11-VOC-SWA
% mdparagraph:P064:begin
A finite \textbf{swā} clause can mean “as; in the same way as”, comparing actions. Use the indicative when both actions are asserted. This differs from the degree adverb swā in swā … ðæt.
% mdparagraph:P064:end

% teaching: NEW-U11-GRAM-u11-verbal-comparison-ma
% teaching: NEW-U11-VOC-MA
% mdparagraph:P065:begin
\textbf{Mā}, “more”, can compare the amount of an action. \textbf{Mā … ðonne} means “more … than”. The second subject remains nominative; a repeated predicate can be understood from the first clause. In a question about what people do, do not turn adverbial mā into an adjective describing the people. Distinguish comparative ðonne, “than”, from temporal ðonne, “when”, and interrogative hwonne, “when?”.
% mdparagraph:P065:end

% teaching: NEW-U11-GRAM-u11-temporal-relative-the
% mdparagraph:P066:begin
After a time noun, \textbf{ðe} can introduce a temporal relative: “the day when …”. It attaches the event to that stated day. It is neither a direct question with hwonne nor a conditional with gif. The ordinary subject-relative construction from Unit 8 also remains available: a demonstrative meaning “those” may be followed by ðe, “who”, with a plural verb for the people concerned.
% mdparagraph:P066:end

\reserve{7\baselineskip}
% mdparagraph:P067:begin
\section*{4. Nouns: complete models and stem changes}
% mdparagraph:P067:end

% teaching: NEW-U11-LEXICAL-DIAGNOSTICS
% mdparagraph:P068:begin
For each new noun learn its meaning, gender and stem together. The diagnostic forms below are genitive singular, dative singular and nominative plural. Where a noun follows a printed model, choose the cell with the required case and number, replace the model stem by the listed stem, and keep the ending. An endingless cell stays endingless, except where an explicit full-headword rule says otherwise. Apply the operation to \textbf{all four cases in both numbers}.
% mdparagraph:P068:end

% mdparagraph:P069:begin
Singular and plural are available in the stated meanings unless a limit is given. The plurals of \textbf{lēad, īsern, regn, frōfor, dēað, fȳr} and \textbf{seolfor} are reference forms: no new count meaning such as separate “rains” is required. \textbf{Sunne} and \textbf{mōna} denote the singular sun and moon here; their plurals are reference. \textbf{Hǣlend} is the proper title and has only the four assigned singular cells.
% mdparagraph:P069:end

\reserve{7\baselineskip}
% mdparagraph:P070:begin
\subsection*{Ordinary masculine and neuter endings}
% mdparagraph:P070:end

% teaching: NEW-U11-VOC-WEALL
% mdparagraph:P071:begin
\textbf{weall}, wall (masculine). Stem \textbf{weall-}. Genitive singular \textbf{wealles}; dative singular \textbf{wealle}; nominative plural \textbf{weallas}. This is the complete ordinary consonant masculine model.
% mdparagraph:P071:end

% mdparagraph:P072:begin
Masculine nominative and accusative singular have no ending. The singular genitive has -es and dative -e; the plural has -as, -a, -um, -as.
% mdparagraph:P072:end

% teaching: NEW-U11-NOM-WEALL
% table: NEW-U11-NOM-WEALL
% mdtable:T01:begin
\par\begingroup
\fontsize{10.5}{13.2}\selectfont
\setlength{\LTleft}{0pt}\setlength{\LTright}{\fill}
\noindent
\begin{tabular}{@{}>{\raggedright\arraybackslash}p{39.46mm}>{\raggedright\arraybackslash}p{39.46mm}>{\raggedright\arraybackslash}p{39.46mm}@{}}
\toprule
\textbf{Case} & \textbf{Singular} & \textbf{Plural} \\
\midrule
Nominative & weall & weallas \\
Genitive & wealles & wealla \\
Dative & wealle & weallum \\
Accusative & weall & weallas \\
\bottomrule
\end{tabular}
\par\endgroup

% mdtable:T01:end
% teaching: NEW-U11-VOC-GIMM
% teaching: NEW-U11-NOM-GIMM
% mdparagraph:P073:begin
\textbf{gimm}, gem; precious stone (masculine). Stem \textbf{gimm-}. Genitive singular \textbf{gimmes}; dative singular \textbf{gimme}; nominative plural \textbf{gimmas}. Use every corresponding cell of \textbf{weall}, replacing weall- by the listed stem.
% mdparagraph:P073:end

% teaching: NEW-U11-VOC-REGN
% teaching: NEW-U11-NOM-REGN
% mdparagraph:P074:begin
\textbf{regn}, rain (masculine). Stem \textbf{regn-}. Genitive singular \textbf{regnes}; dative singular \textbf{regne}; nominative plural \textbf{regnas}. Use every corresponding cell of \textbf{weall}, replacing weall- by the listed stem. The plural is reference only. Cook also gives \textbf{rēn}; regn is the form selected for composition.
% mdparagraph:P074:end

% teaching: NEW-U11-VOC-TORR
% teaching: NEW-U11-NOM-TORR
% mdparagraph:P075:begin
\textbf{torr}, tower (masculine). Stem \textbf{torr-}. Genitive singular \textbf{torres}; dative singular \textbf{torre}; nominative plural \textbf{torras}. Use every corresponding cell of \textbf{weall}, replacing weall- by the listed stem.
% mdparagraph:P075:end

% teaching: NEW-U11-VOC-DEATH
% teaching: NEW-U11-NOM-DEATH
% mdparagraph:P076:begin
\textbf{dēað}, death (masculine). Stem \textbf{dēað-}. Genitive singular \textbf{dēaðes}; dative singular \textbf{dēaðe}; nominative plural \textbf{dēaðas}. Use every corresponding cell of \textbf{weall}, replacing weall- by the listed stem. The plural is reference only.
% mdparagraph:P076:end

% teaching: NEW-U11-VOC-IGLAND
% mdparagraph:P077:begin
\textbf{īgland}, island (neuter). Stem \textbf{īgland-}. Genitive singular \textbf{īglandes}; dative singular \textbf{īglande}; nominative plural \textbf{īgland}. This is the complete long-stem neuter model.
% mdparagraph:P077:end

% mdparagraph:P078:begin
The neuter shares singular -es and -e with the masculine, but nominative and accusative plural have no ending. Plural genitive -a and dative -um remain.
% mdparagraph:P078:end

% teaching: NEW-U11-NOM-IGLAND
% table: NEW-U11-NOM-IGLAND
% mdtable:T02:begin
\par\begingroup
\fontsize{10.5}{13.2}\selectfont
\setlength{\LTleft}{0pt}\setlength{\LTright}{\fill}
\noindent
\begin{tabular}{@{}>{\raggedright\arraybackslash}p{39.46mm}>{\raggedright\arraybackslash}p{39.46mm}>{\raggedright\arraybackslash}p{39.46mm}@{}}
\toprule
\textbf{Case} & \textbf{Singular} & \textbf{Plural} \\
\midrule
Nominative & īgland & īgland \\
Genitive & īglandes & īglanda \\
Dative & īglande & īglandum \\
Accusative & īgland & īgland \\
\bottomrule
\end{tabular}
\par\endgroup

% mdtable:T02:end
% teaching: NEW-U11-VOC-MERESWIN
% teaching: NEW-U11-NOM-MERESWIN
% mdparagraph:P079:begin
\textbf{męreswīn}, dolphin (neuter). Stem \textbf{męreswīn-}. Genitive singular \textbf{męreswīnes}; dative singular \textbf{męreswīne}; nominative plural \textbf{męreswīn}. Use every corresponding cell of \textbf{īgland}, replacing īgland- by the listed stem.
% mdparagraph:P079:end

% teaching: NEW-U11-VOC-LEAD
% teaching: NEW-U11-NOM-LEAD
% mdparagraph:P080:begin
\textbf{lēad}, lead (metal) (neuter). Stem \textbf{lēad-}. Genitive singular \textbf{lēades}; dative singular \textbf{lēade}; nominative plural \textbf{lēad}. Use every corresponding cell of \textbf{īgland}, replacing īgland- by the listed stem. The plural is reference only.
% mdparagraph:P080:end

% teaching: NEW-U11-VOC-ISERN
% teaching: NEW-U11-NOM-ISERN
% mdparagraph:P081:begin
\textbf{īsern}, iron (metal) (neuter). Stem \textbf{īsern-}. Genitive singular \textbf{īsernes}; dative singular \textbf{īserne}; nominative plural \textbf{īsern}. Use every corresponding cell of \textbf{īgland}, replacing īgland- by the listed stem. The plural is reference only. Keep its medial e before endings: the final syllable ends in rn.
% mdparagraph:P081:end

% teaching: NEW-U11-VOC-BEARN
% teaching: NEW-U11-NOM-BEARN
% mdparagraph:P082:begin
\textbf{bearn}, son; child (neuter). Stem \textbf{bearn-}. Genitive singular \textbf{bearnes}; dative singular \textbf{bearne}; nominative plural \textbf{bearn}. Use every corresponding cell of \textbf{īgland}, replacing īgland- by the listed stem. This ordinary “child; son” use does not alter the earlier capitalized title Bearn.
% mdparagraph:P082:end

% teaching: NEW-U11-VOC-FYR
% teaching: NEW-U11-NOM-FYR
% mdparagraph:P083:begin
\textbf{fȳr}, fire (neuter). Stem \textbf{fȳr-}. Genitive singular \textbf{fȳres}; dative singular \textbf{fȳre}; nominative plural \textbf{fȳr}. Use every corresponding cell of \textbf{īgland}, replacing īgland- by the listed stem. The plural is reference only.
% mdparagraph:P083:end

\reserve{7\baselineskip}
% mdparagraph:P084:begin
\subsection*{Vowels that change with number}
% mdparagraph:P084:end

% teaching: NEW-U11-NOMINAL-VOWELS
% mdparagraph:P085:begin
The following vowel changes belong to the named nouns. They cannot be applied to every æ or ea. \textbf{Bæð} keeps æ throughout the singular and has a throughout the plural. \textbf{Geat} keeps ea in the singular and has a in the plural. Their nominative and accusative plurals end in -u.
% mdparagraph:P085:end

% teaching: NEW-U11-VOC-BAETH
% mdparagraph:P086:begin
\textbf{bæð}, bath (neuter). Stem \textbf{bæð-}. Genitive singular \textbf{bæðes}; dative singular \textbf{bæðe}; nominative plural \textbf{baðu}.
% mdparagraph:P086:end

% teaching: NEW-U11-NOM-BAETH
% table: NEW-U11-NOM-BAETH
% mdtable:T03:begin
\par\begingroup
\fontsize{10.5}{13.2}\selectfont
\setlength{\LTleft}{0pt}\setlength{\LTright}{\fill}
\noindent
\begin{tabular}{@{}>{\raggedright\arraybackslash}p{39.46mm}>{\raggedright\arraybackslash}p{39.46mm}>{\raggedright\arraybackslash}p{39.46mm}@{}}
\toprule
\textbf{Case} & \textbf{Singular} & \textbf{Plural} \\
\midrule
Nominative & bæð & baðu \\
Genitive & bæðes & baða \\
Dative & bæðe & baðum \\
Accusative & bæð & baðu \\
\bottomrule
\end{tabular}
\par\endgroup

% mdtable:T03:end
% teaching: NEW-U11-VOC-GEAT
% mdparagraph:P087:begin
\textbf{geat}, gate (neuter). Stem \textbf{geat-}. Genitive singular \textbf{geates}; dative singular \textbf{geate}; nominative plural \textbf{gatu}.
% mdparagraph:P087:end

% teaching: NEW-U11-NOM-GEAT
% table: NEW-U11-NOM-GEAT
% mdtable:T04:begin
\par\begingroup
\fontsize{10.5}{13.2}\selectfont
\setlength{\LTleft}{0pt}\setlength{\LTright}{\fill}
\noindent
\begin{tabular}{@{}>{\raggedright\arraybackslash}p{39.46mm}>{\raggedright\arraybackslash}p{39.46mm}>{\raggedright\arraybackslash}p{39.46mm}@{}}
\toprule
\textbf{Case} & \textbf{Singular} & \textbf{Plural} \\
\midrule
Nominative & geat & gatu \\
Genitive & geates & gata \\
Dative & geate & gatum \\
Accusative & geat & gatu \\
\bottomrule
\end{tabular}
\par\endgroup

% mdtable:T04:end
% mdparagraph:P088:begin
The selected dative plural is \textbf{gatum}. Cook also prints \textbf{geatum} elsewhere; it is an accepted variant, not an extra form required from memory.
% mdparagraph:P088:end

% teaching: NEW-U11-VOC-SITHFAET
% mdparagraph:P089:begin
\textbf{sīðfæt}, journey; passage (masculine). Stem \textbf{sīðfæt-}. Genitive singular \textbf{sīðfætes}; dative singular \textbf{sīðfæte}; nominative plural \textbf{sīðfatas}. The singular uses sīðfæt- and the plural \textbf{sīðfat-}. This noun is masculine, so its nominative and accusative plural take -as.
% mdparagraph:P089:end

% teaching: NEW-U11-NOM-SITHFAET
% table: NEW-U11-NOM-SITHFAET
% mdtable:T05:begin
\par\begingroup
\fontsize{10.5}{13.2}\selectfont
\setlength{\LTleft}{0pt}\setlength{\LTright}{\fill}
\noindent
\begin{tabular}{@{}>{\raggedright\arraybackslash}p{39.46mm}>{\raggedright\arraybackslash}p{39.46mm}>{\raggedright\arraybackslash}p{39.46mm}@{}}
\toprule
\textbf{Case} & \textbf{Singular} & \textbf{Plural} \\
\midrule
Nominative & sīðfæt & sīðfatas \\
Genitive & sīðfætes & sīðfata \\
Dative & sīðfæte & sīðfatum \\
Accusative & sīðfæt & sīðfatas \\
\bottomrule
\end{tabular}
\par\endgroup

% mdtable:T05:end
\reserve{7\baselineskip}
% mdparagraph:P090:begin
\subsection*{Seasons and h-loss}
% mdparagraph:P090:end

% teaching: NEW-U11-NOM-SUMER
% teaching: NEW-U11-VOC-SUMER
% mdparagraph:P091:begin
\textbf{sumer}, summer (masculine). Stem \textbf{sumer-}. Genitive singular \textbf{sumeres}; dative singular \textbf{sumere}; nominative plural \textbf{sumeras}. Keep the medial e before each ending.
% mdparagraph:P091:end

% table: NEW-U11-NOM-SUMER
% mdtable:T06:begin
\par\begingroup
\fontsize{10.5}{13.2}\selectfont
\setlength{\LTleft}{0pt}\setlength{\LTright}{\fill}
\noindent
\begin{tabular}{@{}>{\raggedright\arraybackslash}p{39.46mm}>{\raggedright\arraybackslash}p{39.46mm}>{\raggedright\arraybackslash}p{39.46mm}@{}}
\toprule
\textbf{Case} & \textbf{Singular} & \textbf{Plural} \\
\midrule
Nominative & sumer & sumeras \\
Genitive & sumeres & sumera \\
Dative & sumere & sumerum \\
Accusative & sumer & sumeras \\
\bottomrule
\end{tabular}
\par\endgroup

% mdtable:T06:end
% teaching: NEW-U11-VOC-WINTER
% mdparagraph:P092:begin
\textbf{winter}, winter (masculine). Stem \textbf{wintr-}. Genitive singular \textbf{wintres}; dative singular \textbf{wintre}; nominative plural \textbf{wintras}. The full form winter occurs in nominative and accusative singular; before endings use \textbf{wintr-}.
% mdparagraph:P092:end

% teaching: NEW-U11-NOM-WINTER
% table: NEW-U11-NOM-WINTER
% mdtable:T07:begin
\par\begingroup
\fontsize{10.5}{13.2}\selectfont
\setlength{\LTleft}{0pt}\setlength{\LTright}{\fill}
\noindent
\begin{tabular}{@{}>{\raggedright\arraybackslash}p{39.46mm}>{\raggedright\arraybackslash}p{39.46mm}>{\raggedright\arraybackslash}p{39.46mm}@{}}
\toprule
\textbf{Case} & \textbf{Singular} & \textbf{Plural} \\
\midrule
Nominative & winter & wintras \\
Genitive & wintres & wintra \\
Dative & wintre & wintrum \\
Accusative & winter & wintras \\
\bottomrule
\end{tabular}
\par\endgroup

% mdtable:T07:end
% mdparagraph:P093:begin
Cook records occasional \textbf{sumera} and \textbf{wintra} also as genitive singular, dative singular and nominative plural. They are acceptable alternatives in those cells. Learn the ordinary forms displayed here; the variants need not be recalled independently. The source use of winter for “storm” would be a separate meaning.
% mdparagraph:P093:end

% teaching: NEW-U11-NOM-SEOLH
% teaching: NEW-U11-VOC-SEOLH
% mdparagraph:P094:begin
\textbf{seolh}, seal (masculine). Stem \textbf{sēol-}. Genitive singular \textbf{sēoles}; dative singular \textbf{sēole}; nominative plural \textbf{sēolas}. The bare nominative and accusative singular keep short eo and h. Before an ending, h disappears and the diphthong lengthens: \textbf{sēol-}.
% mdparagraph:P094:end

% table: NEW-U11-NOM-SEOLH
% mdtable:T08:begin
\par\begingroup
\fontsize{10.5}{13.2}\selectfont
\setlength{\LTleft}{0pt}\setlength{\LTright}{\fill}
\noindent
\begin{tabular}{@{}>{\raggedright\arraybackslash}p{39.46mm}>{\raggedright\arraybackslash}p{39.46mm}>{\raggedright\arraybackslash}p{39.46mm}@{}}
\toprule
\textbf{Case} & \textbf{Singular} & \textbf{Plural} \\
\midrule
Nominative & seolh & sēolas \\
Genitive & sēoles & sēola \\
Dative & sēole & sēolum \\
Accusative & seolh & sēolas \\
\bottomrule
\end{tabular}
\par\endgroup

% mdtable:T08:end
% mdparagraph:P095:begin
Both parts of the change matter: sēoles has no h and has long ēo. Cook also records a bare form seol; the table retains seolh.
% mdparagraph:P095:end

\reserve{7\baselineskip}
% mdparagraph:P096:begin
\subsection*{Feminines and internal vowel loss}
% mdparagraph:P096:end

% teaching: NEW-U11-VOC-STRAET
% mdparagraph:P097:begin
\textbf{strǣt}, street (feminine). Stem \textbf{strǣt-}. Genitive singular \textbf{strǣte}; dative singular \textbf{strǣte}; nominative plural \textbf{strǣta}. This is the complete long-stem feminine model.
% mdparagraph:P097:end

% mdparagraph:P098:begin
The nominative singular is endingless. The other singular cases have -e. Nominative, genitive and accusative plural have -a; dative plural has -um.
% mdparagraph:P098:end

% teaching: NEW-U11-NOM-STRAET
% table: NEW-U11-NOM-STRAET
% mdtable:T09:begin
\par\begingroup
\fontsize{10.5}{13.2}\selectfont
\setlength{\LTleft}{0pt}\setlength{\LTright}{\fill}
\noindent
\begin{tabular}{@{}>{\raggedright\arraybackslash}p{39.46mm}>{\raggedright\arraybackslash}p{39.46mm}>{\raggedright\arraybackslash}p{39.46mm}@{}}
\toprule
\textbf{Case} & \textbf{Singular} & \textbf{Plural} \\
\midrule
Nominative & strǣt & strǣta \\
Genitive & strǣte & strǣta \\
Dative & strǣte & strǣtum \\
Accusative & strǣte & strǣta \\
\bottomrule
\end{tabular}
\par\endgroup

% mdtable:T09:end
% teaching: NEW-U11-VOC-GEBYRD
% teaching: NEW-U11-NOM-GEBYRD
% mdparagraph:P099:begin
\textbf{gebyrd}, birth (feminine). Stem \textbf{gebyrd-}. Genitive singular \textbf{gebyrde}; dative singular \textbf{gebyrde}; nominative plural \textbf{gebyrda}. Replace strǣt- by \textbf{gebyrd-} in every cell of strǣt.
% mdparagraph:P099:end

% teaching: NEW-U11-VOC-MILDHEORTNES
% teaching: NEW-U11-NOM-MILDHEORTNES
% mdparagraph:P100:begin
\textbf{mildheortnes}, mercy; mercies (feminine). Stem \textbf{mildheortness-}. Genitive singular \textbf{mildheortnesse}; dative singular \textbf{mildheortnesse}; nominative plural \textbf{mildheortnessa}. Use the endings of \textbf{strǣt} with \textbf{mildheortness-}, but keep the full headword \textbf{mildheortnes} in the endingless nominative singular. The final s is doubled before each ending. Both “mercy” and plural “mercies” are used here.
% mdparagraph:P100:end

% teaching: NEW-U11-VOC-FROFOR
% mdparagraph:P101:begin
\textbf{frōfor}, comfort; consolation (feminine). Stem \textbf{frōfr-}. Genitive singular \textbf{frōfre}; dative singular \textbf{frōfre}; nominative plural \textbf{frōfra}. The nominative singular is the full \textbf{frōfor}; every ending-bearing cell uses \textbf{frōfr-}. Its plural is reference.
% mdparagraph:P101:end

% teaching: NEW-U11-NOM-FROFOR
% table: NEW-U11-NOM-FROFOR
% mdtable:T10:begin
\par\begingroup
\fontsize{10.5}{13.2}\selectfont
\setlength{\LTleft}{0pt}\setlength{\LTright}{\fill}
\noindent
\begin{tabular}{@{}>{\raggedright\arraybackslash}p{39.46mm}>{\raggedright\arraybackslash}p{39.46mm}>{\raggedright\arraybackslash}p{39.46mm}@{}}
\toprule
\textbf{Case} & \textbf{Singular} & \textbf{Plural} \\
\midrule
Nominative & frōfor & frōfra \\
Genitive & frōfre & frōfra \\
Dative & frōfre & frōfrum \\
Accusative & frōfre & frōfra \\
\bottomrule
\end{tabular}
\par\endgroup

% mdtable:T10:end
% teaching: NEW-U11-VOC-CEASTER
% teaching: NEW-U11-NOM-CEASTER
% mdparagraph:P102:begin
\textbf{ceaster}, town; city (feminine). Stem \textbf{ceastr-}. Genitive singular \textbf{ceastre}; dative singular \textbf{ceastre}; nominative plural \textbf{ceastra}. Use \textbf{frōfor} as the full model: replace its ending-bearing stem frōfr- by \textbf{ceastr-}, and its full nominative singular frōfor by \textbf{ceaster}. Both numbers are available.
% mdparagraph:P102:end

% teaching: NEW-U11-VOC-SEOLFOR
% mdparagraph:P103:begin
\textbf{seolfor}, silver (metal) (neuter). Stem \textbf{seolfr-}. Genitive singular \textbf{seolfres}; dative singular \textbf{seolfre}; nominative plural \textbf{seolfor}. Keep the full headword \textbf{seolfor} in all endingless nominative and accusative cells; use \textbf{seolfr-} before endings. The plural is reference.
% mdparagraph:P103:end

% teaching: NEW-U11-NOM-SEOLFOR
% table: NEW-U11-NOM-SEOLFOR
% mdtable:T11:begin
\par\begingroup
\fontsize{10.5}{13.2}\selectfont
\setlength{\LTleft}{0pt}\setlength{\LTright}{\fill}
\noindent
\begin{tabular}{@{}>{\raggedright\arraybackslash}p{39.46mm}>{\raggedright\arraybackslash}p{39.46mm}>{\raggedright\arraybackslash}p{39.46mm}@{}}
\toprule
\textbf{Case} & \textbf{Singular} & \textbf{Plural} \\
\midrule
Nominative & seolfor & seolfor \\
Genitive & seolfres & seolfra \\
Dative & seolfre & seolfrum \\
Accusative & seolfor & seolfor \\
\bottomrule
\end{tabular}
\par\endgroup

% mdtable:T11:end
% mdparagraph:P104:begin
This loss of an unstressed internal vowel is \textbf{syncope}. Compare frōfr-, ceastr-, seolfr- and wintr-. Sumer and īsern retain their vowels. Learn the lexical distinction rather than deleting every internal vowel before an ending.
% mdparagraph:P104:end

\reserve{7\baselineskip}
% mdparagraph:P105:begin
\subsection*{Weak nouns}
% mdparagraph:P105:end

% teaching: NEW-U11-VOC-MEREGREOTE
% mdparagraph:P106:begin
\textbf{męregrēote}, pearl (feminine). Stem \textbf{męregrēot-}. Genitive singular \textbf{męregrēotan}; dative singular \textbf{męregrēotan}; nominative plural \textbf{męregrēotan}. The weak feminine has nominative singular -e, the other singular cases -an, plural nominative and accusative -an, genitive -ena and dative -um.
% mdparagraph:P106:end

% teaching: NEW-U11-NOM-MEREGREOTE
% table: NEW-U11-NOM-MEREGREOTE
% mdtable:T12:begin
\par\begingroup
\fontsize{10.5}{13.2}\selectfont
\setlength{\LTleft}{0pt}\setlength{\LTright}{\fill}
\noindent
\begin{tabular}{@{}>{\raggedright\arraybackslash}p{35.87mm}>{\raggedright\arraybackslash}p{39.46mm}>{\raggedright\arraybackslash}p{43.05mm}@{}}
\toprule
\textbf{Case} & \textbf{Singular} & \textbf{Plural} \\
\midrule
Nominative & męregrēote & męregrēotan \\
Genitive & męregrēotan & męregrēotena \\
Dative & męregrēotan & męregrēotum \\
Accusative & męregrēotan & męregrēotan \\
\bottomrule
\end{tabular}
\par\endgroup

% mdtable:T12:end
% teaching: NEW-U11-VOC-SUNNE
% teaching: NEW-U11-NOM-SUNNE
% mdparagraph:P107:begin
\textbf{sunne}, sun (feminine). Stem \textbf{sunn-}. Genitive singular \textbf{sunnan}; dative singular \textbf{sunnan}; nominative plural \textbf{sunnan}. Use every cell of \textbf{męregrēote}, replacing męregrēot- by \textbf{sunn-}. Thus nominative singular has sunn- + -e; the other singular cases have sunn- + -an. The astronomical singular is for composition; the plural is reference.
% mdparagraph:P107:end

% teaching: NEW-U11-VOC-MONA
% mdparagraph:P108:begin
\textbf{mōna}, moon (masculine). Stem \textbf{mōn-}. Genitive singular \textbf{mōnan}; dative singular \textbf{mōnan}; nominative plural \textbf{mōnan}. The weak masculine differs from the weak feminine in nominative singular -a. Its astronomical singular is for composition; the plural is reference.
% mdparagraph:P108:end

% teaching: NEW-U11-NOM-MONA
% table: NEW-U11-NOM-MONA
% mdtable:T13:begin
\par\begingroup
\fontsize{10.5}{13.2}\selectfont
\setlength{\LTleft}{0pt}\setlength{\LTright}{\fill}
\noindent
\begin{tabular}{@{}>{\raggedright\arraybackslash}p{39.46mm}>{\raggedright\arraybackslash}p{39.46mm}>{\raggedright\arraybackslash}p{39.46mm}@{}}
\toprule
\textbf{Case} & \textbf{Singular} & \textbf{Plural} \\
\midrule
Nominative & mōna & mōnan \\
Genitive & mōnan & mōnena \\
Dative & mōnan & mōnum \\
Accusative & mōnan & mōnan \\
\bottomrule
\end{tabular}
\par\endgroup

% mdtable:T13:end
\reserve{7\baselineskip}
% mdparagraph:P109:begin
\subsection*{Titles and apposition}
% mdparagraph:P109:end

% teaching: NEW-U11-VOC-HAELEND
% teaching: NEW-U11-NOM-HAELEND
% mdparagraph:P110:begin
\textbf{Hǣlend}, the Saviour; Jesus (proper title) (masculine). Stem \textbf{Hǣlend-}. Genitive singular \textbf{Hǣlendes}; dative singular \textbf{Hǣlende}. Use the singular cells of \textbf{weall}, replacing weall- by \textbf{Hǣlend-}. The four title forms are also given together below. There is no assigned proper-title plural.
% mdparagraph:P110:end

% mdtable:T14:begin
\par\begingroup
\fontsize{10.5}{13.2}\selectfont
\setlength{\LTleft}{0pt}\setlength{\LTright}{\fill}
\noindent
\begin{tabular}{@{}>{\raggedright\arraybackslash}p{58.54mm}>{\raggedright\arraybackslash}p{62.64mm}@{}}
\toprule
\textbf{Case} & \textbf{Singular title} \\
\midrule
Nominative & Hǣlend \\
Genitive & Hǣlendes \\
Dative & Hǣlende \\
Accusative & Hǣlend \\
\bottomrule
\end{tabular}
\par\endgroup

% mdtable:T14:end
% teaching: NEW-U11-GRAM-NEW-U11-NOM-APPOSITION
% mdparagraph:P111:begin
Nouns in \textbf{apposition} name the same person or thing and agree in case. In \textbf{ūres Dryhtnes Hǣlendes Crīstes}, “of our Lord Jesus Christ”, the three names or titles are genitive singular; ūres agrees with the masculine singular title. Hǣlend supplies the proper-title equivalent of “Jesus” in this group. In a blessing naming God and the Father, the appositional names remain nominative if they name the subject; a dependent “of …” phrase inside them remains genitive.
% mdparagraph:P111:end

% mdparagraph:P112:begin
Keep the earlier distinction between the case of a whole phrase and a genitive inside it. The composition models normally use the practiced preposed genitive. A supported genitive after its head can also be valid Old English, as the exact reading shows; it is not made ungrammatical by that teaching preference.
% mdparagraph:P112:end

\reserve{7\baselineskip}
% mdparagraph:P113:begin
\section*{5. Adjectives: full models, comparatives and personal uses}
% mdparagraph:P113:end

% teaching: NEW-U11-ADJECTIVE-DIAGNOSTICS
% mdparagraph:P114:begin
Use a strong adjective for the ordinary indefinite or predicative construction and a weak adjective in the taught definite construction after the article. Keep the grammatical gender of the noun, including neuter wīf, “wife; woman”. The tables give all four ordinary cases in all three genders and both numbers. Distinct masculine and neuter singular instrumentals are \textbf{reference}; “as dative” points to the corresponding dative cell and introduces no new use.
% mdparagraph:P114:end

% mdparagraph:P115:begin
The familiar strong endings include singular masculine accusative -ne, feminine genitive/\allowbreak{}dative -re, plural genitive -ra and dative -um. Weak nominative singular has masculine -a and feminine/\allowbreak{}neuter -e; neuter accusative singular also has -e. The remaining singular cases and plural nominative/\allowbreak{}accusative have -an; plural genitive has the selected -ra and dative -um. Stem changes must be applied along with these endings.
% mdparagraph:P115:end

% teaching: NEW-U11-VOC-HAT
% mdparagraph:P116:begin
\textbf{hāt}, hot. Stem \textbf{hāt-}. The long stem has no -u in strong feminine nominative singular or neuter nominative/\allowbreak{}accusative plural.
% mdparagraph:P116:end

% teaching: NEW-U11-ADJ-HAT-STRONG
\reserve{18\baselineskip}
% mdparagraph:P117:begin
\textbf{hāt — strong}
% mdparagraph:P117:end

% table: NEW-U11-ADJ-HAT-STRONG
% mdtable:T15:begin
% course-table-style: 2026-10
% course-panel:T15-P01:begin
\par\begingroup
\fontsize{10.5}{13.2}\selectfont
\setlength{\LTleft}{0pt}\setlength{\LTright}{\fill}
\noindent
\begin{tabular}{@{}>{\raggedright\arraybackslash}p{22.0mm}>{\raggedright\arraybackslash}p{31.188mm}>{\raggedright\arraybackslash}p{31.188mm}>{\raggedright\arraybackslash}p{31.188mm}@{}}
\toprule
\textbf{Case} & \textbf{Masculine singular} & \textbf{Neuter singular} & \textbf{Feminine singular} \\
\midrule
Nominative & hāt & hāt & hāt \\
Genitive & hātes & hātes & hātre \\
Dative & hātum & hātum & hātre \\
Accusative & hātne & hāt & hāte \\
Instrumental & hāte & hāte & as dative \\
\bottomrule
\end{tabular}
\par\endgroup
% course-panel:T15-P01:end
\par\addvspace{5pt}
% course-panel:T15-P02:begin
\par\begingroup
\fontsize{10.5}{13.2}\selectfont
\setlength{\LTleft}{0pt}\setlength{\LTright}{\fill}
\noindent
\begin{tabular}{@{}>{\raggedright\arraybackslash}p{22.0mm}>{\raggedright\arraybackslash}p{31.188mm}>{\raggedright\arraybackslash}p{31.188mm}>{\raggedright\arraybackslash}p{31.188mm}@{}}
\toprule
\textbf{Case} & \textbf{Masculine plural} & \textbf{Neuter plural} & \textbf{Feminine plural} \\
\midrule
Nominative & hāte & hāt & hāta \\
Genitive & hātra & hātra & hātra \\
Dative & hātum & hātum & hātum \\
Accusative & hāte & hāt & hāta \\
Instrumental & as dative & as dative & as dative \\
\bottomrule
\end{tabular}
\par\endgroup
% course-panel:T15-P02:end
% mdtable:T15:end

% teaching: NEW-U11-ADJ-HAT-WEAK
\reserve{18\baselineskip}
% mdparagraph:P118:begin
\textbf{hāt — weak}
% mdparagraph:P118:end

% table: NEW-U11-ADJ-HAT-WEAK
% mdtable:T16:begin
% course-table-style: 2026-10
% course-panel:T16-P01:begin
\par\begingroup
\fontsize{10.5}{13.2}\selectfont
\setlength{\LTleft}{0pt}\setlength{\LTright}{\fill}
\noindent
\begin{tabular}{@{}>{\raggedright\arraybackslash}p{22.0mm}>{\raggedright\arraybackslash}p{31.188mm}>{\raggedright\arraybackslash}p{31.188mm}>{\raggedright\arraybackslash}p{31.188mm}@{}}
\toprule
\textbf{Case} & \textbf{Masculine singular} & \textbf{Neuter singular} & \textbf{Feminine singular} \\
\midrule
Nominative & hāta & hāte & hāte \\
Genitive & hātan & hātan & hātan \\
Dative & hātan & hātan & hātan \\
Accusative & hātan & hāte & hātan \\
Instrumental & hātan & hātan & as dative \\
\bottomrule
\end{tabular}
\par\endgroup
% course-panel:T16-P01:end
\par\addvspace{5pt}
% course-panel:T16-P02:begin
\par\begingroup
\fontsize{10.5}{13.2}\selectfont
\setlength{\LTleft}{0pt}\setlength{\LTright}{\fill}
\noindent
\begin{tabular}{@{}>{\raggedright\arraybackslash}p{22.0mm}>{\raggedright\arraybackslash}p{31.188mm}>{\raggedright\arraybackslash}p{31.188mm}>{\raggedright\arraybackslash}p{31.188mm}@{}}
\toprule
\textbf{Case} & \textbf{Masculine plural} & \textbf{Neuter plural} & \textbf{Feminine plural} \\
\midrule
Nominative & hātan & hātan & hātan \\
Genitive & hātra & hātra & hātra \\
Dative & hātum & hātum & hātum \\
Accusative & hātan & hātan & hātan \\
Instrumental & as dative & as dative & as dative \\
\bottomrule
\end{tabular}
\par\endgroup
% course-panel:T16-P02:end
% mdtable:T16:end

% teaching: NEW-U11-VOC-RIHTWIS
% teaching: NEW-U11-ADJ-RIHTWIS-STRONG
% teaching: NEW-U11-ADJ-RIHTWIS-WEAK
% mdparagraph:P119:begin
\textbf{rihtwīs}, righteous. Stem \textbf{rihtwīs-}. For every strong and weak cell, replace hāt- in the corresponding \textbf{hāt} table by the listed stem, retaining the ending or lack of ending. Instrumental reference limits remain the same.
% mdparagraph:P119:end

% teaching: NEW-U11-VOC-DEAD
% teaching: NEW-U11-ADJ-DEAD-STRONG
% teaching: NEW-U11-ADJ-DEAD-WEAK
% mdparagraph:P120:begin
\textbf{dēad}, dead. Stem \textbf{dēad-}. For every strong and weak cell, replace hāt- in the corresponding \textbf{hāt} table by the listed stem, retaining the ending or lack of ending. Instrumental reference limits remain the same.
% mdparagraph:P120:end

% teaching: NEW-U11-VOC-BLAEC
% mdparagraph:P121:begin
\textbf{blæc}, black. Stem \textbf{blæc-}. In the strong system, use \textbf{blac-} before a vowel ending, but retain \textbf{blæc-} before a consonantal ending or no ending. The short stem takes -u in strong feminine nominative singular and neuter nominative/\allowbreak{}accusative plural. The weak system uses \textbf{blac-} before vowel endings and \textbf{blæc-} before genitive plural -ra.
% mdparagraph:P121:end

% teaching: NEW-U11-ADJ-BLAEC-STRONG
\reserve{18\baselineskip}
% mdparagraph:P122:begin
\textbf{blæc — strong}
% mdparagraph:P122:end

% table: NEW-U11-ADJ-BLAEC-STRONG
% mdtable:T17:begin
% course-table-style: 2026-10
% course-panel:T17-P01:begin
\par\begingroup
\fontsize{10.5}{13.2}\selectfont
\setlength{\LTleft}{0pt}\setlength{\LTright}{\fill}
\noindent
\begin{tabular}{@{}>{\raggedright\arraybackslash}p{22.0mm}>{\raggedright\arraybackslash}p{31.188mm}>{\raggedright\arraybackslash}p{31.188mm}>{\raggedright\arraybackslash}p{31.188mm}@{}}
\toprule
\textbf{Case} & \textbf{Masculine singular} & \textbf{Neuter singular} & \textbf{Feminine singular} \\
\midrule
Nominative & blæc & blæc & blacu \\
Genitive & blaces & blaces & blæcre \\
Dative & blacum & blacum & blæcre \\
Accusative & blæcne & blæc & blace \\
Instrumental & blace & blace & as dative \\
\bottomrule
\end{tabular}
\par\endgroup
% course-panel:T17-P01:end
\par\addvspace{5pt}
% course-panel:T17-P02:begin
\par\begingroup
\fontsize{10.5}{13.2}\selectfont
\setlength{\LTleft}{0pt}\setlength{\LTright}{\fill}
\noindent
\begin{tabular}{@{}>{\raggedright\arraybackslash}p{22.0mm}>{\raggedright\arraybackslash}p{31.188mm}>{\raggedright\arraybackslash}p{31.188mm}>{\raggedright\arraybackslash}p{31.188mm}@{}}
\toprule
\textbf{Case} & \textbf{Masculine plural} & \textbf{Neuter plural} & \textbf{Feminine plural} \\
\midrule
Nominative & blace & blacu & blaca \\
Genitive & blæcra & blæcra & blæcra \\
Dative & blacum & blacum & blacum \\
Accusative & blace & blacu & blaca \\
Instrumental & as dative & as dative & as dative \\
\bottomrule
\end{tabular}
\par\endgroup
% course-panel:T17-P02:end
% mdtable:T17:end

% teaching: NEW-U11-ADJ-BLAEC-WEAK
\reserve{18\baselineskip}
% mdparagraph:P123:begin
\textbf{blæc — weak}
% mdparagraph:P123:end

% table: NEW-U11-ADJ-BLAEC-WEAK
% mdtable:T18:begin
% course-table-style: 2026-10
% course-panel:T18-P01:begin
\par\begingroup
\fontsize{10.5}{13.2}\selectfont
\setlength{\LTleft}{0pt}\setlength{\LTright}{\fill}
\noindent
\begin{tabular}{@{}>{\raggedright\arraybackslash}p{22.0mm}>{\raggedright\arraybackslash}p{31.188mm}>{\raggedright\arraybackslash}p{31.188mm}>{\raggedright\arraybackslash}p{31.188mm}@{}}
\toprule
\textbf{Case} & \textbf{Masculine singular} & \textbf{Neuter singular} & \textbf{Feminine singular} \\
\midrule
Nominative & blaca & blace & blace \\
Genitive & blacan & blacan & blacan \\
Dative & blacan & blacan & blacan \\
Accusative & blacan & blace & blacan \\
Instrumental & blacan & blacan & as dative \\
\bottomrule
\end{tabular}
\par\endgroup
% course-panel:T18-P01:end
\par\addvspace{5pt}
% course-panel:T18-P02:begin
\par\begingroup
\fontsize{10.5}{13.2}\selectfont
\setlength{\LTleft}{0pt}\setlength{\LTright}{\fill}
\noindent
\begin{tabular}{@{}>{\raggedright\arraybackslash}p{22.0mm}>{\raggedright\arraybackslash}p{31.188mm}>{\raggedright\arraybackslash}p{31.188mm}>{\raggedright\arraybackslash}p{31.188mm}@{}}
\toprule
\textbf{Case} & \textbf{Masculine plural} & \textbf{Neuter plural} & \textbf{Feminine plural} \\
\midrule
Nominative & blacan & blacan & blacan \\
Genitive & blæcra & blæcra & blæcra \\
Dative & blacum & blacum & blacum \\
Accusative & blacan & blacan & blacan \\
Instrumental & as dative & as dative & as dative \\
\bottomrule
\end{tabular}
\par\endgroup
% course-panel:T18-P02:end
% mdtable:T18:end

% mdparagraph:P124:begin
Thus strong masculine accusative blæcne retains æ before -ne; weak masculine nominative blaca has a before -a. The weak genitive plural \textbf{blæcra} uses \textbf{blæc-} before -ra, the ending shared with the strong declension.
% mdparagraph:P124:end

% teaching: NEW-U11-VOC-FAEGER
% mdparagraph:P125:begin
\textbf{fæger}, beautiful; fair. Stem \textbf{fæger-}. The selected disyllabic model retains its internal e. Strong feminine nominative singular and neuter nominative/\allowbreak{}accusative plural have -u.
% mdparagraph:P125:end

% teaching: NEW-U11-ADJ-FAEGER-STRONG
\reserve{18\baselineskip}
% mdparagraph:P126:begin
\textbf{fæger — strong}
% mdparagraph:P126:end

% table: NEW-U11-ADJ-FAEGER-STRONG
% mdtable:T19:begin
% course-table-style: 2026-10
% course-panel:T19-P01:begin
\par\begingroup
\fontsize{10.5}{13.2}\selectfont
\setlength{\LTleft}{0pt}\setlength{\LTright}{\fill}
\noindent
\begin{tabular}{@{}>{\raggedright\arraybackslash}p{22.0mm}>{\raggedright\arraybackslash}p{31.188mm}>{\raggedright\arraybackslash}p{31.188mm}>{\raggedright\arraybackslash}p{31.188mm}@{}}
\toprule
\textbf{Case} & \textbf{Masculine singular} & \textbf{Neuter singular} & \textbf{Feminine singular} \\
\midrule
Nominative & fæger & fæger & fægeru \\
Genitive & fægeres & fægeres & fægerre \\
Dative & fægerum & fægerum & fægerre \\
Accusative & fægerne & fæger & fægere \\
Instrumental & fægere & fægere & as dative \\
\bottomrule
\end{tabular}
\par\endgroup
% course-panel:T19-P01:end
\par\addvspace{5pt}
% course-panel:T19-P02:begin
\par\begingroup
\fontsize{10.5}{13.2}\selectfont
\setlength{\LTleft}{0pt}\setlength{\LTright}{\fill}
\noindent
\begin{tabular}{@{}>{\raggedright\arraybackslash}p{22.0mm}>{\raggedright\arraybackslash}p{31.188mm}>{\raggedright\arraybackslash}p{31.188mm}>{\raggedright\arraybackslash}p{31.188mm}@{}}
\toprule
\textbf{Case} & \textbf{Masculine plural} & \textbf{Neuter plural} & \textbf{Feminine plural} \\
\midrule
Nominative & fægere & fægeru & fægera \\
Genitive & fægerra & fægerra & fægerra \\
Dative & fægerum & fægerum & fægerum \\
Accusative & fægere & fægeru & fægera \\
Instrumental & as dative & as dative & as dative \\
\bottomrule
\end{tabular}
\par\endgroup
% course-panel:T19-P02:end
% mdtable:T19:end

% teaching: NEW-U11-ADJ-FAEGER-WEAK
\reserve{18\baselineskip}
% mdparagraph:P127:begin
\textbf{fæger — weak}
% mdparagraph:P127:end

% table: NEW-U11-ADJ-FAEGER-WEAK
% mdtable:T20:begin
% course-table-style: 2026-10
% course-panel:T20-P01:begin
\par\begingroup
\fontsize{10.5}{13.2}\selectfont
\setlength{\LTleft}{0pt}\setlength{\LTright}{\fill}
\noindent
\begin{tabular}{@{}>{\raggedright\arraybackslash}p{22.0mm}>{\raggedright\arraybackslash}p{31.188mm}>{\raggedright\arraybackslash}p{31.188mm}>{\raggedright\arraybackslash}p{31.188mm}@{}}
\toprule
\textbf{Case} & \textbf{Masculine singular} & \textbf{Neuter singular} & \textbf{Feminine singular} \\
\midrule
Nominative & fægera & fægere & fægere \\
Genitive & fægeran & fægeran & fægeran \\
Dative & fægeran & fægeran & fægeran \\
Accusative & fægeran & fægere & fægeran \\
Instrumental & fægeran & fægeran & as dative \\
\bottomrule
\end{tabular}
\par\endgroup
% course-panel:T20-P01:end
\par\addvspace{5pt}
% course-panel:T20-P02:begin
\par\begingroup
\fontsize{10.5}{13.2}\selectfont
\setlength{\LTleft}{0pt}\setlength{\LTright}{\fill}
\noindent
\begin{tabular}{@{}>{\raggedright\arraybackslash}p{22.0mm}>{\raggedright\arraybackslash}p{31.188mm}>{\raggedright\arraybackslash}p{31.188mm}>{\raggedright\arraybackslash}p{31.188mm}@{}}
\toprule
\textbf{Case} & \textbf{Masculine plural} & \textbf{Neuter plural} & \textbf{Feminine plural} \\
\midrule
Nominative & fægeran & fægeran & fægeran \\
Genitive & fægerra & fægerra & fægerra \\
Dative & fægerum & fægerum & fægerum \\
Accusative & fægeran & fægeran & fægeran \\
Instrumental & as dative & as dative & as dative \\
\bottomrule
\end{tabular}
\par\endgroup
% course-panel:T20-P02:end
% mdtable:T20:end

% teaching: NEW-U11-VOC-GESYNDIG
% teaching: NEW-U11-ADJ-GESYNDIG-STRONG
% teaching: NEW-U11-ADJ-GESYNDIG-WEAK
% mdparagraph:P128:begin
\textbf{gesyndig}, favorable; fair (of winds). Stem \textbf{gesyndig-}. Replace fæger- by the listed stem in \textbf{every cell of both fæger tables}. Keep the internal vowel, including before -u, -re and -ra; the reference instrumentals follow the same operation.
% mdparagraph:P128:end

% teaching: NEW-U11-VOC-THURSTIG
% teaching: NEW-U11-ADJ-THURSTIG-STRONG
% teaching: NEW-U11-ADJ-THURSTIG-WEAK
% mdparagraph:P129:begin
\textbf{ðurstig}, thirsty. Stem \textbf{ðurstig-}. Replace fæger- by the listed stem in \textbf{every cell of both fæger tables}. Keep the internal vowel, including before -u, -re and -ra; the reference instrumentals follow the same operation.
% mdparagraph:P129:end

% mdparagraph:P130:begin
Cook also allows internal-vowel loss in such disyllabic adjectives, and occasional endingless feminine singular and neuter plural forms. The tables choose the retained-vowel series consistently. The exact \textbf{gesyndge}, “favorable”, in the sea passage is a supplied syncopated form beside the selected \textbf{gesyndige}. Do not silently replace one spelling by the other.
% mdparagraph:P130:end

% teaching: NEW-U11-VOC-SMYLTE
% mdparagraph:P131:begin
\textbf{smylte}, calm; smooth. Stem \textbf{smylt-}. The endingless masculine and neuter singular use the full headword \textbf{smylte}. Elsewhere replace final -e by the required ending on \textbf{smylt-}. For example, feminine nominative singular has -u and dative plural -um.
% mdparagraph:P131:end

% teaching: NEW-U11-ADJ-SMYLTE-STRONG
\reserve{18\baselineskip}
% mdparagraph:P132:begin
\textbf{smylte — strong}
% mdparagraph:P132:end

% table: NEW-U11-ADJ-SMYLTE-STRONG
% mdtable:T21:begin
% course-table-style: 2026-10
% course-panel:T21-P01:begin
\par\begingroup
\fontsize{10.5}{13.2}\selectfont
\setlength{\LTleft}{0pt}\setlength{\LTright}{\fill}
\noindent
\begin{tabular}{@{}>{\raggedright\arraybackslash}p{22.0mm}>{\raggedright\arraybackslash}p{31.188mm}>{\raggedright\arraybackslash}p{31.188mm}>{\raggedright\arraybackslash}p{31.188mm}@{}}
\toprule
\textbf{Case} & \textbf{Masculine singular} & \textbf{Neuter singular} & \textbf{Feminine singular} \\
\midrule
Nominative & smylte & smylte & smyltu \\
Genitive & smyltes & smyltes & smyltre \\
Dative & smyltum & smyltum & smyltre \\
Accusative & smyltne & smylte & smylte \\
Instrumental & smylte & smylte & as dative \\
\bottomrule
\end{tabular}
\par\endgroup
% course-panel:T21-P01:end
\par\addvspace{5pt}
% course-panel:T21-P02:begin
\par\begingroup
\fontsize{10.5}{13.2}\selectfont
\setlength{\LTleft}{0pt}\setlength{\LTright}{\fill}
\noindent
\begin{tabular}{@{}>{\raggedright\arraybackslash}p{22.0mm}>{\raggedright\arraybackslash}p{31.188mm}>{\raggedright\arraybackslash}p{31.188mm}>{\raggedright\arraybackslash}p{31.188mm}@{}}
\toprule
\textbf{Case} & \textbf{Masculine plural} & \textbf{Neuter plural} & \textbf{Feminine plural} \\
\midrule
Nominative & smylte & smyltu & smylta \\
Genitive & smyltra & smyltra & smyltra \\
Dative & smyltum & smyltum & smyltum \\
Accusative & smylte & smyltu & smylta \\
Instrumental & as dative & as dative & as dative \\
\bottomrule
\end{tabular}
\par\endgroup
% course-panel:T21-P02:end
% mdtable:T21:end

% teaching: NEW-U11-ADJ-SMYLTE-WEAK
\reserve{18\baselineskip}
% mdparagraph:P133:begin
\textbf{smylte — weak}
% mdparagraph:P133:end

% table: NEW-U11-ADJ-SMYLTE-WEAK
% mdtable:T22:begin
% course-table-style: 2026-10
% course-panel:T22-P01:begin
\par\begingroup
\fontsize{10.5}{13.2}\selectfont
\setlength{\LTleft}{0pt}\setlength{\LTright}{\fill}
\noindent
\begin{tabular}{@{}>{\raggedright\arraybackslash}p{22.0mm}>{\raggedright\arraybackslash}p{31.188mm}>{\raggedright\arraybackslash}p{31.188mm}>{\raggedright\arraybackslash}p{31.188mm}@{}}
\toprule
\textbf{Case} & \textbf{Masculine singular} & \textbf{Neuter singular} & \textbf{Feminine singular} \\
\midrule
Nominative & smylta & smylte & smylte \\
Genitive & smyltan & smyltan & smyltan \\
Dative & smyltan & smyltan & smyltan \\
Accusative & smyltan & smylte & smyltan \\
Instrumental & smyltan & smyltan & as dative \\
\bottomrule
\end{tabular}
\par\endgroup
% course-panel:T22-P01:end
\par\addvspace{5pt}
% course-panel:T22-P02:begin
\par\begingroup
\fontsize{10.5}{13.2}\selectfont
\setlength{\LTleft}{0pt}\setlength{\LTright}{\fill}
\noindent
\begin{tabular}{@{}>{\raggedright\arraybackslash}p{22.0mm}>{\raggedright\arraybackslash}p{31.188mm}>{\raggedright\arraybackslash}p{31.188mm}>{\raggedright\arraybackslash}p{31.188mm}@{}}
\toprule
\textbf{Case} & \textbf{Masculine plural} & \textbf{Neuter plural} & \textbf{Feminine plural} \\
\midrule
Nominative & smyltan & smyltan & smyltan \\
Genitive & smyltra & smyltra & smyltra \\
Dative & smyltum & smyltum & smyltum \\
Accusative & smyltan & smyltan & smyltan \\
Instrumental & as dative & as dative & as dative \\
\bottomrule
\end{tabular}
\par\endgroup
% course-panel:T22-P02:end
% mdtable:T22:end

% teaching: NEW-U11-VOC-STILLE
% mdparagraph:P134:begin
\textbf{stille}, still; quiet. Stem \textbf{still-}. Use full \textbf{stille} in endingless masculine and neuter singular cells, and \textbf{still-} before endings. Simplify ll to l before consonantal -ne, -re and -ra: \textbf{stilne, stilre, stilra}. Retain ll before vowel endings: \textbf{stille, stilles, stillum}. The weak genitive plural is also stilra.
% mdparagraph:P134:end

% teaching: NEW-U11-ADJ-STILLE-STRONG
\reserve{18\baselineskip}
% mdparagraph:P135:begin
\textbf{stille — strong}
% mdparagraph:P135:end

% table: NEW-U11-ADJ-STILLE-STRONG
% mdtable:T23:begin
% course-table-style: 2026-10
% course-panel:T23-P01:begin
\par\begingroup
\fontsize{10.5}{13.2}\selectfont
\setlength{\LTleft}{0pt}\setlength{\LTright}{\fill}
\noindent
\begin{tabular}{@{}>{\raggedright\arraybackslash}p{22.0mm}>{\raggedright\arraybackslash}p{31.188mm}>{\raggedright\arraybackslash}p{31.188mm}>{\raggedright\arraybackslash}p{31.188mm}@{}}
\toprule
\textbf{Case} & \textbf{Masculine singular} & \textbf{Neuter singular} & \textbf{Feminine singular} \\
\midrule
Nominative & stille & stille & stillu \\
Genitive & stilles & stilles & stilre \\
Dative & stillum & stillum & stilre \\
Accusative & stilne & stille & stille \\
Instrumental & stille & stille & as dative \\
\bottomrule
\end{tabular}
\par\endgroup
% course-panel:T23-P01:end
\par\addvspace{5pt}
% course-panel:T23-P02:begin
\par\begingroup
\fontsize{10.5}{13.2}\selectfont
\setlength{\LTleft}{0pt}\setlength{\LTright}{\fill}
\noindent
\begin{tabular}{@{}>{\raggedright\arraybackslash}p{22.0mm}>{\raggedright\arraybackslash}p{31.188mm}>{\raggedright\arraybackslash}p{31.188mm}>{\raggedright\arraybackslash}p{31.188mm}@{}}
\toprule
\textbf{Case} & \textbf{Masculine plural} & \textbf{Neuter plural} & \textbf{Feminine plural} \\
\midrule
Nominative & stille & stillu & stilla \\
Genitive & stilra & stilra & stilra \\
Dative & stillum & stillum & stillum \\
Accusative & stille & stillu & stilla \\
Instrumental & as dative & as dative & as dative \\
\bottomrule
\end{tabular}
\par\endgroup
% course-panel:T23-P02:end
% mdtable:T23:end

% teaching: NEW-U11-ADJ-STILLE-WEAK
\reserve{18\baselineskip}
% mdparagraph:P136:begin
\textbf{stille — weak}
% mdparagraph:P136:end

% table: NEW-U11-ADJ-STILLE-WEAK
% mdtable:T24:begin
% course-table-style: 2026-10
% course-panel:T24-P01:begin
\par\begingroup
\fontsize{10.5}{13.2}\selectfont
\setlength{\LTleft}{0pt}\setlength{\LTright}{\fill}
\noindent
\begin{tabular}{@{}>{\raggedright\arraybackslash}p{22.0mm}>{\raggedright\arraybackslash}p{31.188mm}>{\raggedright\arraybackslash}p{31.188mm}>{\raggedright\arraybackslash}p{31.188mm}@{}}
\toprule
\textbf{Case} & \textbf{Masculine singular} & \textbf{Neuter singular} & \textbf{Feminine singular} \\
\midrule
Nominative & stilla & stille & stille \\
Genitive & stillan & stillan & stillan \\
Dative & stillan & stillan & stillan \\
Accusative & stillan & stille & stillan \\
Instrumental & stillan & stillan & as dative \\
\bottomrule
\end{tabular}
\par\endgroup
% course-panel:T24-P01:end
\par\addvspace{5pt}
% course-panel:T24-P02:begin
\par\begingroup
\fontsize{10.5}{13.2}\selectfont
\setlength{\LTleft}{0pt}\setlength{\LTright}{\fill}
\noindent
\begin{tabular}{@{}>{\raggedright\arraybackslash}p{22.0mm}>{\raggedright\arraybackslash}p{31.188mm}>{\raggedright\arraybackslash}p{31.188mm}>{\raggedright\arraybackslash}p{31.188mm}@{}}
\toprule
\textbf{Case} & \textbf{Masculine plural} & \textbf{Neuter plural} & \textbf{Feminine plural} \\
\midrule
Nominative & stillan & stillan & stillan \\
Genitive & stilra & stilra & stilra \\
Dative & stillum & stillum & stillum \\
Accusative & stillan & stillan & stillan \\
Instrumental & as dative & as dative & as dative \\
\bottomrule
\end{tabular}
\par\endgroup
% course-panel:T24-P02:end
% mdtable:T24:end

% mdparagraph:P137:begin
The literal stille tables show why substituting still- mechanically into smylte without consonant reduction would fail. Quiet or still is an ordinary adjective predicate; it does not call for a passive participle.
% mdparagraph:P137:end

\reserve{7\baselineskip}
% mdparagraph:P138:begin
\subsection*{Comparatives}
% mdparagraph:P138:end

% teaching: NEW-U11-VOC-DEORWIERTH-COMP
% mdparagraph:P139:begin
\textbf{dēorwierðra}, more precious. Stem \textbf{dēorwierðr-}. From dēorwierðe, remove final e and form the comparative stem \textbf{dēorwierðr-}. Comparatives remain weak even as predicates.
% mdparagraph:P139:end

% mdparagraph:P140:begin
Attach the weak endings to the comparative stem. Before genitive plural -ra, the resulting -ðrr- group simplifies to -ðr-: \textbf{dēorwierðra}. Thus masculine nominative singular and genitive plural happen to share a spelling. The instrumental row merely refers to the dative forms.
% mdparagraph:P140:end

% teaching: NEW-U11-ADJ-DEORWIERTH-COMP-WEAK
\reserve{18\baselineskip}
% mdparagraph:P141:begin
\textbf{dēorwierðra — weak}
% mdparagraph:P141:end

% table: NEW-U11-ADJ-DEORWIERTH-COMP-WEAK
% mdtable:T25:begin
% course-table-style: 2026-10
% course-panel:T25-P01:begin
\par\begingroup
\fontsize{10.5}{13.2}\selectfont
\setlength{\LTleft}{0pt}\setlength{\LTright}{\fill}
\noindent
\begin{tabular}{@{}>{\raggedright\arraybackslash}p{22.0mm}>{\raggedright\arraybackslash}p{31.188mm}>{\raggedright\arraybackslash}p{31.188mm}>{\raggedright\arraybackslash}p{31.188mm}@{}}
\toprule
\textbf{Case} & \textbf{Masculine singular} & \textbf{Neuter singular} & \textbf{Feminine singular} \\
\midrule
Nominative & dēorwierðra & dēorwierðre & dēorwierðre \\
Genitive & dēorwierðran & dēorwierðran & dēorwierðran \\
Dative & dēorwierðran & dēorwierðran & dēorwierðran \\
Accusative & dēorwierðran & dēorwierðre & dēorwierðran \\
Instrumental & as dative & as dative & as dative \\
\bottomrule
\end{tabular}
\par\endgroup
% course-panel:T25-P01:end
\par\addvspace{5pt}
% course-panel:T25-P02:begin
\par\begingroup
\fontsize{10.5}{13.2}\selectfont
\setlength{\LTleft}{0pt}\setlength{\LTright}{\fill}
\noindent
\begin{tabular}{@{}>{\raggedright\arraybackslash}p{22.0mm}>{\raggedright\arraybackslash}p{31.188mm}>{\raggedright\arraybackslash}p{31.188mm}>{\raggedright\arraybackslash}p{31.188mm}@{}}
\toprule
\textbf{Case} & \textbf{Masculine plural} & \textbf{Neuter plural} & \textbf{Feminine plural} \\
\midrule
Nominative & dēorwierðran & dēorwierðran & dēorwierðran \\
Genitive & dēorwierðra & dēorwierðra & dēorwierðra \\
Dative & dēorwierðrum & dēorwierðrum & dēorwierðrum \\
Accusative & dēorwierðran & dēorwierðran & dēorwierðran \\
Instrumental & as dative & as dative & as dative \\
\bottomrule
\end{tabular}
\par\endgroup
% course-panel:T25-P02:end
% mdtable:T25:end

% teaching: NEW-U11-VOC-FAEGER-COMP
% teaching: NEW-U11-ADJ-FAEGER-COMP-WEAK
% mdparagraph:P142:begin
\textbf{fægerra}, more beautiful; fairer. Stem \textbf{fægerr-}. Use the entire \textbf{dēorwierðra} table, replacing comparative stem dēorwierðr- by \textbf{fægerr-} and keeping the same weak endings. Before genitive plural -ra, simplify the resulting consonant group to \textbf{fægerra}. The nominative singular is \textbf{fægerra, fægerre, fægerre}; the dative plural is \textbf{fægerrum} in all genders.
% mdparagraph:P142:end

% mdparagraph:P143:begin
A comparative predicate agrees with its own subject: the standard after ðonne has its own gender and case. Gold and silver are neuter; righteousness and night are feminine. Do not make the predicate agree with the noun after “than”.
% mdparagraph:P143:end

\reserve{7\baselineskip}
% mdparagraph:P144:begin
\subsection*{Adjectives used without an expressed noun}
% mdparagraph:P144:end

% teaching: NEW-U11-GRAM-u11-bare-strong-personal-adjectives
% teaching: NEW-U11-VOC-CWIC
% mdparagraph:P145:begin
The bare \textbf{strong genitive plurals dēadra} and \textbf{cwicra} can mean “of the dead” and “of the living”, referring to people. No article is expressed in this selected construction. Cwic continues the already learned adjective “alive”; English “living” does not require a new participle or progressive tense.
% mdparagraph:P145:end

% teaching: NEW-U11-GRAM-u11-articled-righteous-persons
% mdparagraph:P146:begin
With an article, \textbf{rihtwīs} and \textbf{unrihtwīs} can name righteous and unrighteous people. Use the appropriate \textbf{weak} adjective ending. The whole personal group receives the case required by its clause or preposition; after the taught ofer in the exercise, use accusative. This extends the earlier personal-adjective construction to these meanings.
% mdparagraph:P146:end

% teaching: NEW-U11-GRAM-u11-eall-substantive
% teaching: NEW-U11-VOC-EALL
% mdparagraph:P147:begin
\textbf{Eall} can stand alone: neuter singular \textbf{eall}, “everything”, or masculine plural \textbf{ealle}, “all people”. State which meaning is intended when English simply says “all”. In “over all”, the preposition governs accusative; the two meanings keep their distinct numbers.
% mdparagraph:P147:end

% teaching: NEW-U11-GRAM-u11-vocative-personal-adjective
% mdparagraph:P148:begin
The earlier superlative of lēof may also stand alone in singular direct address: \textbf{lēofosta}, “dearest”, to a male person; \textbf{lēofoste} to a female person. These are weak forms. The speaker’s choice of addressee determines the gender; an English prompt without that information admits either defensible choice.
% mdparagraph:P148:end

\reserve{7\baselineskip}
% mdparagraph:P149:begin
\section*{6. Verbs and their complements}
% mdparagraph:P149:end

% mdparagraph:P150:begin
The seven finite tables below give the complete \textbf{present and preterite indicative}, all three persons in both numbers. They do not assign new subjunctive or imperative systems for these verbs. The purpose and other subjunctive exercises use verbs whose full subjunctives have already been learned. Learn each new verb’s case construction together with its forms.
% mdparagraph:P150:end

% teaching: NEW-U11-VOC-SCETHTHAN
% teaching: NEW-U11-GRAM-u11-sceththan-indicative
% teaching: NEW-U11-GRAM-u11-sceththan-dative
% mdparagraph:P151:begin
\textbf{Scęððan}, “hurt; injure”, takes the affected person or thing in the \textbf{dative}. English may make it a direct object, but Old English need not. The selected weak class I system has scęðð- before -e and -að, single ð in \textbf{scęðest, scęðeð}, and an uncontracted dental preterite \textbf{scęð-ede}. Preterite second singular adds -st to the first-singular form; the plural replaces final -e by -on.
% mdparagraph:P151:end

% teaching: NEW-U11-PAR-sceththan-IND
% table: NEW-U11-PAR-sceththan-IND
% mdtable:T26:begin
% course-table-style: 2026-10
% course-panel:T26-P01:begin
\par\begingroup
\fontsize{10.5}{13.2}\selectfont
\setlength{\LTleft}{0pt}\setlength{\LTright}{\fill}
\noindent
\begin{tabular}{@{}>{\raggedright\arraybackslash}p{39.459mm}>{\raggedright\arraybackslash}p{39.459mm}>{\raggedright\arraybackslash}p{39.459mm}@{}}
\toprule
\textbf{Person and number} & \textbf{Present} & \textbf{Preterite} \\
\midrule
1 sg. & scęððe & scęðede \\
2 sg. & scęðest & scęðedest \\
3 sg. & scęðeð & scęðede \\
1–3 pl. & scęððað & scęðedon \\
\bottomrule
\end{tabular}
\par\endgroup
% course-panel:T26-P01:end
% mdtable:T26:end
% mdparagraph:P152:begin
Learn this consistent weak system. Contracted present variants and the separate strong preterite scōd are outside the required series.
% mdparagraph:P152:end

% teaching: NEW-U11-VOC-GEWEORTHIAN
% teaching: NEW-U11-GRAM-u11-geweorthian-indicative
% mdparagraph:P153:begin
\textbf{Geweorðian}, “distinguish; adorn”, is weak class II. It takes the \textbf{accusative} thing adorned or distinguished; \textbf{mid + dative} can name its equipment or adornment. A passive can leave the agent unexpressed. Keep the selected spelling geweorðian.
% mdparagraph:P153:end

% mdparagraph:P154:begin
The class II present has -ie, -ast, -að in the singular and -iað in all three plural persons. The preterite has -ode, -odest, -ode, then -odon in each plural person. The vowel before the dental is o.
% mdparagraph:P154:end

% teaching: NEW-U11-PAR-geweorthian-IND
% table: NEW-U11-PAR-geweorthian-IND
% mdtable:T27:begin
% course-table-style: 2026-10
% course-panel:T27-P01:begin
\par\begingroup
\fontsize{10.5}{13.2}\selectfont
\setlength{\LTleft}{0pt}\setlength{\LTright}{\fill}
\noindent
\begin{tabular}{@{}>{\raggedright\arraybackslash}p{39.459mm}>{\raggedright\arraybackslash}p{39.459mm}>{\raggedright\arraybackslash}p{39.459mm}@{}}
\toprule
\textbf{Person and number} & \textbf{Present} & \textbf{Preterite} \\
\midrule
1 sg. & geweorðie & geweorðode \\
2 sg. & geweorðast & geweorðodest \\
3 sg. & geweorðað & geweorðode \\
1–3 pl. & geweorðiað & geweorðodon \\
\bottomrule
\end{tabular}
\par\endgroup
% course-panel:T27-P01:end
% mdtable:T27:end
% teaching: NEW-U11-VOC-SCEAWIAN
% teaching: NEW-U11-GRAM-u11-sceawian-indicative
% teaching: NEW-U11-GRAM-u11-sceawian-willan
% mdparagraph:P155:begin
\textbf{Scēawian}, “watch; observe”, uses the same class II endings on \textbf{scēaw-}. The person or event watched, including a journey, is \textbf{accusative}. With willan expressing a wish, use plain \textbf{scēawian} as the infinitive.
% mdparagraph:P155:end

% teaching: NEW-U11-PAR-sceawian-IND
% table: NEW-U11-PAR-sceawian-IND
% mdtable:T28:begin
% course-table-style: 2026-10
% course-panel:T28-P01:begin
\par\begingroup
\fontsize{10.5}{13.2}\selectfont
\setlength{\LTleft}{0pt}\setlength{\LTright}{\fill}
\noindent
\begin{tabular}{@{}>{\raggedright\arraybackslash}p{39.459mm}>{\raggedright\arraybackslash}p{39.459mm}>{\raggedright\arraybackslash}p{39.459mm}@{}}
\toprule
\textbf{Person and number} & \textbf{Present} & \textbf{Preterite} \\
\midrule
1 sg. & scēawie & scēawode \\
2 sg. & scēawast & scēawodest \\
3 sg. & scēawað & scēawode \\
1–3 pl. & scēawiað & scēawodon \\
\bottomrule
\end{tabular}
\par\endgroup
% course-panel:T28-P01:end
% mdtable:T28:end
% teaching: NEW-U11-PAR-sceawian-INF
% table: NEW-U11-PAR-sceawian-INF
% mdtable:T29:begin
\par\begingroup
\fontsize{10.5}{13.2}\selectfont
\setlength{\LTleft}{0pt}\setlength{\LTright}{\fill}
\noindent
\begin{tabular}{@{}>{\raggedright\arraybackslash}p{71.72mm}>{\raggedright\arraybackslash}p{49.46mm}@{}}
\toprule
\textbf{Form} & \textbf{Infinitive} \\
\midrule
Plain infinitive & scēawian \\
Inflected infinitive & scēawianne \\
\bottomrule
\end{tabular}
\par\endgroup

% mdtable:T29:end
% mdparagraph:P156:begin
The plain infinitive is for composition after willan. \textbf{Scēawianne} is the inflected infinitive, shown for \textbf{reference}; do not substitute it after willan or attach tō there. The two infinitive shapes do not themselves determine the actor of a finite purpose clause.
% mdparagraph:P156:end

% teaching: NEW-U11-VOC-FEOHTAN
% teaching: NEW-U11-VOC-WITH
% teaching: NEW-U11-GRAM-u11-feohtan-indicative
% teaching: NEW-U11-GRAM-u11-feohtan-opponent
% teaching: NEW-U11-GRAM-u11-with-accusative
% mdparagraph:P157:begin
\textbf{Feohtan}, “fight; contend”, is strong class III. The present has eo except for selected \textbf{fiehtst} and \textbf{fieht} in the second and third singular. The preterite has \textbf{feaht} in first and third singular, but \textbf{fuht-} in second singular and all plural persons: fuhte, fuhton. There is no added dental preterite suffix.
% mdparagraph:P157:end

% teaching: NEW-U11-PAR-feohtan-IND
% table: NEW-U11-PAR-feohtan-IND
% mdtable:T30:begin
% course-table-style: 2026-10
% course-panel:T30-P01:begin
\par\begingroup
\fontsize{10.5}{13.2}\selectfont
\setlength{\LTleft}{0pt}\setlength{\LTright}{\fill}
\noindent
\begin{tabular}{@{}>{\raggedright\arraybackslash}p{39.459mm}>{\raggedright\arraybackslash}p{39.459mm}>{\raggedright\arraybackslash}p{39.459mm}@{}}
\toprule
\textbf{Person and number} & \textbf{Present} & \textbf{Preterite} \\
\midrule
1 sg. & feohte & feaht \\
2 sg. & fiehtst & fuhte \\
3 sg. & fieht & feaht \\
1–3 pl. & feohtað & fuhton \\
\bottomrule
\end{tabular}
\par\endgroup
% course-panel:T30-P01:end
% mdtable:T30:end
% mdparagraph:P158:begin
The opponent follows \textbf{wið + accusative}, “against; with” in the sense of fighting someone. This is opposition, not accompaniment with mid.
% mdparagraph:P158:end

% teaching: NEW-U11-VOC-GENEALAECAN
% teaching: NEW-U11-GRAM-u11-genealaecan-indicative
% teaching: NEW-U11-GRAM-u11-genealaecan-dative
% mdparagraph:P159:begin
\textbf{Genēalǣcan}, “approach; draw near”, takes a \textbf{dative} person, place or event. A time can draw near an event in the same construction. The long stem is \textbf{genēalǣc-}; present second and third singular directly add -st and -ð. In the selected weak preterite the voiceless c is followed by dental t: \textbf{genēalǣc-te}.
% mdparagraph:P159:end

% teaching: NEW-U11-PAR-genealaecan-IND
% table: NEW-U11-PAR-genealaecan-IND
% mdtable:T31:begin
% course-table-style: 2026-10
% course-panel:T31-P01:begin
\par\begingroup
\fontsize{10.5}{13.2}\selectfont
\setlength{\LTleft}{0pt}\setlength{\LTright}{\fill}
\noindent
\begin{tabular}{@{}>{\raggedright\arraybackslash}p{39.459mm}>{\raggedright\arraybackslash}p{39.459mm}>{\raggedright\arraybackslash}p{39.459mm}@{}}
\toprule
\textbf{Person and number} & \textbf{Present} & \textbf{Preterite} \\
\midrule
1 sg. & genēalǣce & genēalǣcte \\
2 sg. & genēalǣcst & genēalǣctest \\
3 sg. & genēalǣcð & genēalǣcte \\
1–3 pl. & genēalǣcað & genēalǣcton \\
\bottomrule
\end{tabular}
\par\endgroup
% course-panel:T31-P01:end
% mdtable:T31:end
% mdparagraph:P160:begin
The second singular preterite has -test and the plural -ton. Attested \textbf{genēalǣhte} is an accepted alternative; the complete selected series follows Cook’s -cte formation and is the one to learn here.
% mdparagraph:P160:end

% teaching: NEW-U11-VOC-HYDAN
% teaching: NEW-U11-GRAM-u11-hydan-indicative
% mdparagraph:P161:begin
\textbf{Hȳdan}, “hide; conceal”, takes an \textbf{accusative} object. Its weak class I dental stem has contracted present \textbf{hȳtst, hȳt}, but hȳd- before -e and -að. The preterite adds the dental suffix to hȳd-, giving \textbf{hȳdde}; keep both d letters. The resulting hidden state may be expressed by the passive construction below.
% mdparagraph:P161:end

% teaching: NEW-U11-PAR-hydan-IND
% table: NEW-U11-PAR-hydan-IND
% mdtable:T32:begin
% course-table-style: 2026-10
% course-panel:T32-P01:begin
\par\begingroup
\fontsize{10.5}{13.2}\selectfont
\setlength{\LTleft}{0pt}\setlength{\LTright}{\fill}
\noindent
\begin{tabular}{@{}>{\raggedright\arraybackslash}p{39.459mm}>{\raggedright\arraybackslash}p{39.459mm}>{\raggedright\arraybackslash}p{39.459mm}@{}}
\toprule
\textbf{Person and number} & \textbf{Present} & \textbf{Preterite} \\
\midrule
1 sg. & hȳde & hȳdde \\
2 sg. & hȳtst & hȳddest \\
3 sg. & hȳt & hȳdde \\
1–3 pl. & hȳdað & hȳddon \\
\bottomrule
\end{tabular}
\par\endgroup
% course-panel:T32-P01:end
% mdtable:T32:end
% teaching: NEW-U11-VOC-BRINGAN
% teaching: NEW-U11-GRAM-u11-bringan-indicative
% teaching: NEW-U11-GRAM-u11-bringan-destination
% mdparagraph:P162:begin
\textbf{Bringan}, “bring”, takes the person or thing brought in the \textbf{accusative}, with \textbf{tō + dative} for its destination, which may itself be a person. Present second singular is \textbf{bringst}, third singular \textbf{bringð}. The selected preterite is weak \textbf{brōhte}, with changed stem \textbf{brōh-} plus dental t. The vowel change does not by itself make the preterite strong.
% mdparagraph:P162:end

% teaching: NEW-U11-PAR-bringan-IND
% table: NEW-U11-PAR-bringan-IND
% mdtable:T33:begin
% course-table-style: 2026-10
% course-panel:T33-P01:begin
\par\begingroup
\fontsize{10.5}{13.2}\selectfont
\setlength{\LTleft}{0pt}\setlength{\LTright}{\fill}
\noindent
\begin{tabular}{@{}>{\raggedright\arraybackslash}p{39.459mm}>{\raggedright\arraybackslash}p{39.459mm}>{\raggedright\arraybackslash}p{39.459mm}@{}}
\toprule
\textbf{Person and number} & \textbf{Present} & \textbf{Preterite} \\
\midrule
1 sg. & bringe & brōhte \\
2 sg. & bringst & brōhtest \\
3 sg. & bringð & brōhte \\
1–3 pl. & bringað & brōhton \\
\bottomrule
\end{tabular}
\par\endgroup
% course-panel:T33-P01:end
% mdtable:T33:end
% mdparagraph:P163:begin
The second singular preterite adds -st to brōhte; all plural persons have brōhton. A separate strong variant is not required.
% mdparagraph:P163:end

\reserve{7\baselineskip}
% mdparagraph:P164:begin
\subsection*{Participles and passive predicates}
% mdparagraph:P164:end

% teaching: NEW-U11-GRAM-u11-passive-predicate-bases
% teaching: NEW-U11-GRAM-u11-fon-catch-accusative
% mdparagraph:P165:begin
\textbf{Fōn}, “catch”, now takes the caught animal, person or thing in the \textbf{accusative}. In the passive, that entity is the subject. The earlier passive with wesan can use a bare past participle, including with a plural subject. Use that construction with the new bases below. Do not invent plural adjective endings from the base alone.
% mdparagraph:P165:end

% mdparagraph:P166:begin
The present-participle ending is \textbf{-nde}: the ordinary linking vowel is e, class II has -iende, and contracted fōn has \textbf{fōnde}. The strong past participles have -en with their own vowel stems: \textbf{fang-en, fund-en}. The weak past participles have the dental ending \textbf{-ed} or class II \textbf{-od}. These are the two base forms of each verb, not full participial adjective declensions.
% mdparagraph:P166:end

% teaching: NEW-U11-GRAM-u11-fon-participle-bases
\reserve{7\baselineskip}
% mdparagraph:P167:begin
\textbf{fōn.}
% mdparagraph:P167:end

% teaching: NEW-U11-PAR-fon-PART
% table: NEW-U11-PAR-fon-PART
% mdtable:T34:begin
\par\begingroup
\fontsize{10.5}{13.2}\selectfont
\setlength{\LTleft}{0pt}\setlength{\LTright}{\fill}
\noindent
\begin{tabular}{@{}>{\raggedright\arraybackslash}p{69.84mm}>{\raggedright\arraybackslash}p{51.35mm}@{}}
\toprule
\textbf{Form} & \textbf{Base} \\
\midrule
Present participle & fōnde \\
Past participle & fangen \\
\bottomrule
\end{tabular}
\par\endgroup

% mdtable:T34:end
% teaching: NEW-U11-GRAM-u11-findan-participle-bases
\reserve{7\baselineskip}
% mdparagraph:P168:begin
\textbf{findan.}
% mdparagraph:P168:end

% teaching: NEW-U11-PAR-findan-PART
% table: NEW-U11-PAR-findan-PART
% mdtable:T35:begin
\par\begingroup
\fontsize{10.5}{13.2}\selectfont
\setlength{\LTleft}{0pt}\setlength{\LTright}{\fill}
\noindent
\begin{tabular}{@{}>{\raggedright\arraybackslash}p{69.84mm}>{\raggedright\arraybackslash}p{51.35mm}@{}}
\toprule
\textbf{Form} & \textbf{Base} \\
\midrule
Present participle & findende \\
Past participle & funden \\
\bottomrule
\end{tabular}
\par\endgroup

% mdtable:T35:end
% teaching: NEW-U11-GRAM-u11-fedan-participle-bases
\reserve{7\baselineskip}
% mdparagraph:P169:begin
\textbf{fēdan.}
% mdparagraph:P169:end

% teaching: NEW-U11-PAR-fedan-PART
% table: NEW-U11-PAR-fedan-PART
% mdtable:T36:begin
\par\begingroup
\fontsize{10.5}{13.2}\selectfont
\setlength{\LTleft}{0pt}\setlength{\LTright}{\fill}
\noindent
\begin{tabular}{@{}>{\raggedright\arraybackslash}p{69.84mm}>{\raggedright\arraybackslash}p{51.35mm}@{}}
\toprule
\textbf{Form} & \textbf{Base} \\
\midrule
Present participle & fēdende \\
Past participle & fēded \\
\bottomrule
\end{tabular}
\par\endgroup

% mdtable:T36:end
% teaching: NEW-U11-GRAM-u11-geweorthian-participle-bases
\reserve{7\baselineskip}
% mdparagraph:P170:begin
\textbf{geweorðian.}
% mdparagraph:P170:end

% teaching: NEW-U11-PAR-geweorthian-PART
% table: NEW-U11-PAR-geweorthian-PART
% mdtable:T37:begin
\par\begingroup
\fontsize{10.5}{13.2}\selectfont
\setlength{\LTleft}{0pt}\setlength{\LTright}{\fill}
\noindent
\begin{tabular}{@{}>{\raggedright\arraybackslash}p{64.38mm}>{\raggedright\arraybackslash}p{56.81mm}@{}}
\toprule
\textbf{Form} & \textbf{Base} \\
\midrule
Present participle & geweorðiende \\
Past participle & geweorðod \\
\bottomrule
\end{tabular}
\par\endgroup

% mdtable:T37:end
% teaching: NEW-U11-GRAM-u11-hydan-participle-bases
\reserve{7\baselineskip}
% mdparagraph:P171:begin
\textbf{hȳdan.}
% mdparagraph:P171:end

% teaching: NEW-U11-PAR-hydan-PART
% table: NEW-U11-PAR-hydan-PART
% mdtable:T38:begin
\par\begingroup
\fontsize{10.5}{13.2}\selectfont
\setlength{\LTleft}{0pt}\setlength{\LTright}{\fill}
\noindent
\begin{tabular}{@{}>{\raggedright\arraybackslash}p{69.84mm}>{\raggedright\arraybackslash}p{51.35mm}@{}}
\toprule
\textbf{Form} & \textbf{Base} \\
\midrule
Present participle & hȳdende \\
Past participle & hȳded \\
\bottomrule
\end{tabular}
\par\endgroup

% mdtable:T38:end
% mdparagraph:P172:begin
The past participles fēded and hȳded also have contracted forms fēd and hȳd. These are acceptable variants; the tables retain fēded and hȳded as the selected bases.
% mdparagraph:P172:end

% mdparagraph:P173:begin
Thus the plural subject still controls a plural auxiliary, while the chosen predicate can remain bare. Present participle forms are learned morphologically here; their display does not create a general progressive tense. Nor do these pairs introduce a new perfect tense or all possible prefixed variants. The exact readings supply their own additional realizations.
% mdparagraph:P173:end

\reserve{7\baselineskip}
% mdparagraph:P174:begin
\subsection*{Departure commands}
% mdparagraph:P174:end

% teaching: NEW-U11-GRAM-u11-gewitan-imperative
% mdparagraph:P175:begin
The earlier \textbf{gewītan}, “depart; go away”, now has both direct commands. Singular \textbf{gewīt} uses the bare long-vowel stem; plural \textbf{gewītað} adds -að. Retain the earlier departure construction \textbf{fram + dative} for the person or place left.
% mdparagraph:P175:end

% teaching: NEW-U11-PAR-gewitan-IMP
% table: NEW-U11-PAR-gewitan-IMP
% mdtable:T39:begin
% course-table-style: 2026-10
% course-panel:T39-P01:begin
\par\begingroup
\fontsize{10.5}{13.2}\selectfont
\setlength{\LTleft}{0pt}\setlength{\LTright}{\fill}
\noindent
\begin{tabular}{@{}>{\raggedright\arraybackslash}p{63.342mm}>{\raggedright\arraybackslash}p{57.846mm}@{}}
\toprule
\textbf{Person and number} & \textbf{Command} \\
\midrule
2 sg. & gewīt \\
2 pl. & gewītað \\
\bottomrule
\end{tabular}
\par\endgroup
% course-panel:T39-P01:end
% mdtable:T39:end
\reserve{7\baselineskip}
% mdparagraph:P176:begin
\section*{7. Additional meanings and constructions}
% mdparagraph:P176:end

% teaching: NEW-U11-VOC-HIERAN
% teaching: NEW-U11-GRAM-u11-hieran-obey-dative
% mdparagraph:P177:begin
\textbf{Hīeran}, besides “hear”, can mean \textbf{“obey” with a dative} person. Coordinated subjects require a plural verb. Do not infer an accusative from the English object “him”.
% mdparagraph:P177:end

% teaching: NEW-U11-VOC-DON
% teaching: NEW-U11-VOC-ON
% teaching: NEW-U11-GRAM-u11-don-placement
% mdparagraph:P178:begin
\textbf{Dōn}, “put; place”, takes an \textbf{accusative} object. For placement into or onto something, use \textbf{on + accusative} as the goal. This differs from the already taught on + dative for location. In a negative command, retain the taught ne with the appropriate singular or plural command.
% mdparagraph:P178:end

% teaching: NEW-U11-VOC-THURH
% teaching: NEW-U11-GRAM-u11-thurh-accusative
% mdparagraph:P179:begin
\textbf{Þurh + accusative}, “through”, gives the spatial path. Learn the þ spelling for composition. The exact sea reading has \textbf{ðurh}, with ð, and retains it in both its path phrase and the supplied expression ðurh eall, “entirely”.
% mdparagraph:P179:end

% teaching: NEW-U11-VOC-RISAN
% teaching: NEW-U11-VOC-UP
% teaching: NEW-U11-VOC-ONGEAN
% teaching: NEW-U11-GRAM-u11-risan-up-against
% mdparagraph:P180:begin
The familiar \textbf{rīsan}, “rise”, combines with \textbf{ūp}, “up”. For rising against an opponent or adverse power, add \textbf{ongēan + accusative}, “against”. Ūp itself gives no opponent. Keep the already learned unprefixed rīsan forms; this phrase does not require a new verb ārīsan.
% mdparagraph:P180:end

% teaching: NEW-U11-VOC-BUTAN
% teaching: NEW-U11-VOC-WIF
% teaching: NEW-U11-GRAM-u11-butan-dative
% mdparagraph:P181:begin
\textbf{Būtan + dative}, “besides; apart from”, excludes a group from the number just stated. In this use \textbf{wīf} means “woman”, as well as its earlier “wife”; it remains grammatically neuter. Keep the case governed by būtan when the phrase names women and children.
% mdparagraph:P181:end

% teaching: NEW-U11-VOC-BIDDAN
% teaching: NEW-U11-GRAM-u11-biddan-person
% mdparagraph:P182:begin
In the parenthetic \textbf{ic ðē bidde}, “I beseech you; please”, biddan takes the person entreated in the \textbf{accusative}. The same surface ðē can be dative elsewhere, but the verb determines its case here. This bounded request does not introduce every possible complement of biddan.
% mdparagraph:P182:end

% teaching: NEW-U11-VOC-DOM
% teaching: NEW-U11-VOC-OFT
% mdparagraph:P183:begin
\textbf{Dōm} can mean “doom; fate” in addition to the earlier “judgment”. Keep its familiar masculine forms. \textbf{Oft} means “often”; like hēr and mā it is an invariant adverb.
% mdparagraph:P183:end

% teaching: NEW-U11-VOC-NALAES
% teaching: NEW-U11-VOC-HOLUNGA
% teaching: NEW-U11-GRAM-u11-nalaes-negation
% mdparagraph:P184:begin
Learn \textbf{nālæs hōlunga}, “not at all in vain”, as the selected emphatic phrase: nālæs, “not at all”, denies hōlunga, “in vain”. It says that the action was not fruitless. Do not mechanically add another ne. This phrase does not teach the general system of negative concord.
% mdparagraph:P184:end

% mdparagraph:P185:begin
The ordinary senses and forms from earlier units remain in force: helpan governs a dative person; sęllan takes an accusative thing and dative recipient; clothing takes an accusative person; possessive his and hiera name the possessor without agreeing as an adjective. Familiar meanings do not erase the particular new case constructions just learned.
% mdparagraph:P185:end

\reserve{7\baselineskip}
% mdparagraph:P186:begin
\section*{8. Practice: forms}
% mdparagraph:P186:end

% mdparagraph:P187:begin
Give the requested forms, retaining their vowel lengths and spelling. Explain the relevant formation rather than merely naming a class. Reference variants are accepted where noted, but are not additional memorization targets.
% mdparagraph:P187:end

\reserve{5\baselineskip}
\par\begin{minipage}{\linewidth}\raggedright
% mdparagraph:P188:begin
\textbf{F01.}
% mdparagraph:P188:end

% item: NEW-U11-F01
% mdparagraph:P189:begin
Give the complete nominative, genitive, dative and accusative singular and plural of bæð and geat. Then give the nominative plural and dative plural of sīðfæt. Explain where each lexical vowel changes.
% mdparagraph:P189:end

\end{minipage}\par\addvspace{4pt}

\par\begin{minipage}{\linewidth}\raggedright
% mdparagraph:P190:begin
\textbf{F02.}
% mdparagraph:P190:end

% item: NEW-U11-F02
% mdparagraph:P191:begin
Give the complete nominative, genitive, dative and accusative singular and plural of seolh. Explain which cells retain h and which have the long-vowel stem.
% mdparagraph:P191:end

\end{minipage}\par\addvspace{4pt}

\par\begin{minipage}{\linewidth}\raggedright
% mdparagraph:P192:begin
\textbf{F03.}
% mdparagraph:P192:end

% item: NEW-U11-F03
% mdparagraph:P193:begin
Give the complete nominative, genitive, dative and accusative singular and plural of sumer and winter, using the selected tables. Explain why the stems before endings differ. The noted -a variants are acceptable but are not required from memory.
% mdparagraph:P193:end

\end{minipage}\par\addvspace{4pt}

\par\begin{minipage}{\linewidth}\raggedright
% mdparagraph:P194:begin
\textbf{F04.}
% mdparagraph:P194:end

% item: NEW-U11-F04
% mdparagraph:P195:begin
Give these forms: blæc, strong masculine accusative singular, strong feminine nominative singular, and weak masculine nominative singular; stille, strong masculine accusative singular; fægerra, nominative singular in each of the three genders; dēorwierðra, dative plural in all genders. Explain the stem and weak-ending choices.
% mdparagraph:P195:end

\end{minipage}\par\addvspace{4pt}

\par\begin{minipage}{\linewidth}\raggedright
% mdparagraph:P196:begin
\textbf{F05.}
% mdparagraph:P196:end

% item: NEW-U11-F05
% mdparagraph:P197:begin
Give the complete present and preterite indicative of scęððan and genēalǣcan: first, second and third person in both numbers for each tense. Explain the selected stem and dental-preterite formation.
% mdparagraph:P197:end

\end{minipage}\par\addvspace{4pt}

\par\begin{minipage}{\linewidth}\raggedright
% mdparagraph:P198:begin
\textbf{F06.}
% mdparagraph:P198:end

% item: NEW-U11-F06
% mdparagraph:P199:begin
For feohtan and bringan give the present indicative second and third singular, the preterite first and second singular, and the preterite first plural. Explain which preterite belongs to a strong vowel series and which contains a dental suffix.
% mdparagraph:P199:end

\end{minipage}\par\addvspace{4pt}

\par\begin{minipage}{\linewidth}\raggedright
% mdparagraph:P200:begin
\textbf{F07.}
% mdparagraph:P200:end

% item: NEW-U11-F07
% mdparagraph:P201:begin
Give the base present and past participles of fōn, findan, fēdan, geweorðian and hȳdan. Identify the present-participle ending and the strong or weak past-participle ending in each pair.
% mdparagraph:P201:end

\end{minipage}\par\addvspace{4pt}

\par\begin{minipage}{\linewidth}\raggedright
% mdparagraph:P202:begin
\textbf{F08.}
% mdparagraph:P202:end

% item: NEW-U11-F08
% mdparagraph:P203:begin
Give the singular and plural commands of gewītan. Then give hȳdan preterite indicative third singular and third plural, and scēawian and geweorðian present indicative third singular and third plural. Explain what changes when one addressee or subject becomes several.
% mdparagraph:P203:end

\end{minipage}\par\addvspace{4pt}

\reserve{7\baselineskip}
% mdparagraph:P204:begin
\section*{9. Practice: clauses and meaning}
% mdparagraph:P204:end

% mdparagraph:P205:begin
Translate each printed Old English input and answer the questions attached to it. Where two or three inputs are given, treat all of them. The stated situation is part of the task.
% mdparagraph:P205:end

\reserve{5\baselineskip}
\par\begin{minipage}{\linewidth}\raggedright
% mdparagraph:P206:begin
\textbf{A01.}
% mdparagraph:P206:end

% item: NEW-U11-A01
% mdparagraph:P207:begin
Translate. Mark the subordinate and main clauses, identify both finite verbs, and explain the two occurrences of ðā.
% mdparagraph:P207:end

% mdparagraph:P208:begin
Ðā se munuc cōm, ðā sealde se mann him wæter.
% mdparagraph:P208:end

\end{minipage}\par\addvspace{4pt}

\par\begin{minipage}{\linewidth}\raggedright
% mdparagraph:P209:begin
\textbf{A02.}
% mdparagraph:P209:end

% item: NEW-U11-A02
% mdparagraph:P210:begin
The speaker describes what happens on repeated visits. Translate; identify the clause introduced by the first ðonne and the role of the second.
% mdparagraph:P210:end

% mdparagraph:P211:begin
Ðonne gē cumað, ðonne helpe ic ēow.
% mdparagraph:P211:end

\end{minipage}\par\addvspace{4pt}

\par\begin{minipage}{\linewidth}\raggedright
% mdparagraph:P212:begin
\textbf{A03.}
% mdparagraph:P212:end

% item: NEW-U11-A03
% mdparagraph:P213:begin
Translate both statements. Explain the relationship between the actions and whether either statement by itself says that one action caused the other.
% mdparagraph:P213:end

% mdparagraph:P214:begin
(a) Ðā-hwīle-ðe hē ǣt, ic bær wæter.
% mdparagraph:P214:end

% mdparagraph:P215:begin
(b) Mid-ðȳ hē ǣt, ic bær wæter.
% mdparagraph:P215:end

\end{minipage}\par\addvspace{4pt}

\par\begin{minipage}{\linewidth}\raggedright
% mdparagraph:P216:begin
\textbf{A04.}
% mdparagraph:P216:end

% item: NEW-U11-A04
% mdparagraph:P217:begin
Both instructions concern a coming that has not yet occurred. Translate and identify the mood of cume; explain why the English future setting does not require a new Old English future ending.
% mdparagraph:P217:end

% mdparagraph:P218:begin
(a) Ne far, ǣr ic cume.
% mdparagraph:P218:end

% mdparagraph:P219:begin
(b) Wes hēr, oð-ðæt ic cume.
% mdparagraph:P219:end

\end{minipage}\par\addvspace{4pt}

\par\begin{minipage}{\linewidth}\raggedright
% mdparagraph:P220:begin
\textbf{A05.}
% mdparagraph:P220:end

% item: NEW-U11-A05
% mdparagraph:P221:begin
Translate. State which clause gives the reason and which gives the action being explained; identify the finite form in the reason clause.
% mdparagraph:P221:end

% mdparagraph:P222:begin
Ic helpe ðē, for-ðon-ðe ðū eart untrum.
% mdparagraph:P222:end

\end{minipage}\par\addvspace{4pt}

\par\begin{minipage}{\linewidth}\raggedright
% mdparagraph:P223:begin
\textbf{A06.}
% mdparagraph:P223:end

% item: NEW-U11-A06
% mdparagraph:P224:begin
Translate each sentence and classify its ðæt clause as content, intended purpose or asserted result. Parse the finite verb inside that clause.
% mdparagraph:P224:end

% mdparagraph:P225:begin
(a) Ic wāt ðæt se munuc cymð.
% mdparagraph:P225:end

% mdparagraph:P226:begin
(b) Se munuc cymð, ðæt hē mē helpe.
% mdparagraph:P226:end

% mdparagraph:P227:begin
(c) Se munuc is swā gōd, ðæt hē mē hilpð.
% mdparagraph:P227:end

\end{minipage}\par\addvspace{4pt}

\par\begin{minipage}{\linewidth}\raggedright
% mdparagraph:P228:begin
\textbf{A07.}
% mdparagraph:P228:end

% item: NEW-U11-A07
% mdparagraph:P229:begin
Translate. Identify the outcome the man intends to prevent; explain whether the sentence states that the animal actually took the food.
% mdparagraph:P229:end

% mdparagraph:P230:begin
Se mann hēold ðæt dēor, ðȳ-lǣs-ðe hit ðone męte nōme.
% mdparagraph:P230:end

\end{minipage}\par\addvspace{4pt}

\par\begin{minipage}{\linewidth}\raggedright
% mdparagraph:P231:begin
\textbf{A08.}
% mdparagraph:P231:end

% item: NEW-U11-A08
% mdparagraph:P232:begin
Translate and parse sīe. The speaker accepts that the man is sick: can the selected subjunctive still be used, and what relation does the clause express?
% mdparagraph:P232:end

% mdparagraph:P233:begin
Ðēah hē sīe untrum, hē hilpð mē.
% mdparagraph:P233:end

\end{minipage}\par\addvspace{4pt}

\par\begin{minipage}{\linewidth}\raggedright
% mdparagraph:P234:begin
\textbf{A09.}
% mdparagraph:P234:end

% item: NEW-U11-A09
% mdparagraph:P235:begin
In (a), his presence is an open question. In (b), he is absent now. Translate both, identify their finite moods, and state the time referred to by wǣre and hulpe.
% mdparagraph:P235:end

% mdparagraph:P236:begin
(a) Gif hē is hēr, hē hilpð mē.
% mdparagraph:P236:end

% mdparagraph:P237:begin
(b) Gif hē wǣre hēr, hē hulpe mē.
% mdparagraph:P237:end

\end{minipage}\par\addvspace{4pt}

\par\begin{minipage}{\linewidth}\raggedright
% mdparagraph:P238:begin
\textbf{A10.}
% mdparagraph:P238:end

% item: NEW-U11-A10
% mdparagraph:P239:begin
Translate each sentence. Distinguish interrogative hwonne, temporal ðonne, and comparative ðonne; do not infer the function from the spelling alone.
% mdparagraph:P239:end

% mdparagraph:P240:begin
(a) Hwonne cymð se munuc?
% mdparagraph:P240:end

% mdparagraph:P241:begin
(b) Ðonne se munuc cymð, ic helpe him.
% mdparagraph:P241:end

% mdparagraph:P242:begin
(c) Se munuc is ieldra ðonne se fiscere.
% mdparagraph:P242:end

\end{minipage}\par\addvspace{4pt}

\par\begin{minipage}{\linewidth}\raggedright
% mdparagraph:P243:begin
\textbf{A11.}
% mdparagraph:P243:end

% item: NEW-U11-A11
% mdparagraph:P244:begin
Translate. Who is understood to do the helping in each sentence? Explain why a finite purpose clause can make a different subject explicit.
% mdparagraph:P244:end

% mdparagraph:P245:begin
(a) Se fiscere cōm tō helpanne ðǣm męnn.
% mdparagraph:P245:end

% mdparagraph:P246:begin
(b) Se fiscere cōm, ðæt se munuc ðǣm męnn hulpe.
% mdparagraph:P246:end

\end{minipage}\par\addvspace{4pt}

\par\begin{minipage}{\linewidth}\raggedright
% mdparagraph:P247:begin
\textbf{A12.}
% mdparagraph:P247:end

% item: NEW-U11-A12
% mdparagraph:P248:begin
Translate. Identify the antecedent of ðe and explain why this is neither a direct when-question nor a free-standing conditional.
% mdparagraph:P248:end

% mdparagraph:P249:begin
Hit is se dæg ðe hē cōm.
% mdparagraph:P249:end

\end{minipage}\par\addvspace{4pt}

\reserve{7\baselineskip}
% mdparagraph:P250:begin
\section*{10. Controlled composition}
% mdparagraph:P250:end

% mdparagraph:P251:begin
Translate into Old English using the stated construction. A grammatical alternative may still miss a requested target; distinguish those two judgments. Preserve the people, number, time, purpose and other relations specified in the English.
% mdparagraph:P251:end

\reserve{5\baselineskip}
\par\begin{minipage}{\linewidth}\raggedright
% mdparagraph:P252:begin
\textbf{T01.}
% mdparagraph:P252:end

% item: NEW-U11-T01
% mdparagraph:P253:begin
Translate using past ðā … ðā: When the father came, the son gave him water.
% mdparagraph:P253:end

\end{minipage}\par\addvspace{4pt}

\par\begin{minipage}{\linewidth}\raggedright
% mdparagraph:P254:begin
\textbf{T02.}
% mdparagraph:P254:end

% item: NEW-U11-T02
% mdparagraph:P255:begin
Translate as a recurring event, using ðonne … ðonne: When the fishermen come, the king gives them food.
% mdparagraph:P255:end

\end{minipage}\par\addvspace{4pt}

\par\begin{minipage}{\linewidth}\raggedright
% mdparagraph:P256:begin
\textbf{T03.}
% mdparagraph:P256:end

% item: NEW-U11-T03
% mdparagraph:P257:begin
Translate with mid-ðȳ: While the fishermen sat in the house, the son carried water.
% mdparagraph:P257:end

\end{minipage}\par\addvspace{4pt}

\par\begin{minipage}{\linewidth}\raggedright
% mdparagraph:P258:begin
\textbf{T04.}
% mdparagraph:P258:end

% item: NEW-U11-T04
% mdparagraph:P259:begin
Address one person. The speaker has not yet arrived. Translate with ǣr: Do not go before I come.
% mdparagraph:P259:end

\end{minipage}\par\addvspace{4pt}

\par\begin{minipage}{\linewidth}\raggedright
% mdparagraph:P260:begin
\textbf{T05.}
% mdparagraph:P260:end

% item: NEW-U11-T05
% mdparagraph:P261:begin
Address several people. The fishermen have not yet arrived. Translate with oð-ðæt: Be in the house until the fishermen come.
% mdparagraph:P261:end

\end{minipage}\par\addvspace{4pt}

\par\begin{minipage}{\linewidth}\raggedright
% mdparagraph:P262:begin
\textbf{T06.}
% mdparagraph:P262:end

% item: NEW-U11-T06
% mdparagraph:P263:begin
Translate with for-ðon-ðe and the adjective ðurstig: The monk gives the man water because the man is thirsty.
% mdparagraph:P263:end

\end{minipage}\par\addvspace{4pt}

\par\begin{minipage}{\linewidth}\raggedright
% mdparagraph:P264:begin
\textbf{T07.}
% mdparagraph:P264:end

% item: NEW-U11-T07
% mdparagraph:P265:begin
Translate with a finite ðæt purpose clause: The fishermen came so that the monk might help them.
% mdparagraph:P265:end

\end{minipage}\par\addvspace{4pt}

\par\begin{minipage}{\linewidth}\raggedright
% mdparagraph:P266:begin
\textbf{T08.}
% mdparagraph:P266:end

% item: NEW-U11-T08
% mdparagraph:P267:begin
Translate the result as an actual past event, using swā … ðæt: The men were so hungry that they ate the food.
% mdparagraph:P267:end

\end{minipage}\par\addvspace{4pt}

\par\begin{minipage}{\linewidth}\raggedright
% mdparagraph:P268:begin
\textbf{T09.}
% mdparagraph:P268:end

% item: NEW-U11-T09
% mdparagraph:P269:begin
Translate with ðȳ-lǣs-ðe: The fishermen held the animals, lest the animals take our food.
% mdparagraph:P269:end

\end{minipage}\par\addvspace{4pt}

\par\begin{minipage}{\linewidth}\raggedright
% mdparagraph:P270:begin
\textbf{T10.}
% mdparagraph:P270:end

% item: NEW-U11-T10
% mdparagraph:P271:begin
The monk was in fact sick. Translate using the selected concessive pattern: Although the monk was sick, he helped the man.
% mdparagraph:P271:end

\end{minipage}\par\addvspace{4pt}

\par\begin{minipage}{\linewidth}\raggedright
% mdparagraph:P272:begin
\textbf{T11.}
% mdparagraph:P272:end

% item: NEW-U11-T11
% mdparagraph:P273:begin
The king is absent now. Translate the contrary-present condition: If the king were here, he would help us.
% mdparagraph:P273:end

\end{minipage}\par\addvspace{4pt}

\reserve{7\baselineskip}
% mdparagraph:P274:begin
\section*{11. Cook IV.16–17 and V.18}
% mdparagraph:P274:end

% mdparagraph:P275:begin
Translate the three complete exercise sets. Their original wording, grammatical references and parenthetical directions are retained. \textbf{Om.} means “omit”: it applies to the indicated English word, not to the meaning of the sentence. \textbf{P.p.} means “past participle”. Material labelled local help is supplied for use; do not treat it as an additional recall test.
% mdparagraph:P275:end

\reserve{7\baselineskip}
% mdparagraph:P276:begin
\subsection*{IV.16}
% mdparagraph:P276:end

% mdparagraph:P277:begin
For the unspecified “your” command in IV.16.16, state whether you address one person or several and keep the command and possessive consistent. For “all” in IV.16.13, identify whether you mean everything or all people. Use the participle bases from §6 in the passive items.
% mdparagraph:P277:end

\reserve{5\baselineskip}
\par\begin{minipage}{\linewidth}\raggedright
% mdparagraph:P278:begin
\textbf{IV.16.1.}
% mdparagraph:P278:end

% item: EX-IV-16-01
% mdparagraph:P279:begin
Gold is more precious (64) than silver.
% mdparagraph:P279:end

\end{minipage}\par\addvspace{4pt}

\par\begin{minipage}{\linewidth}\raggedright
% mdparagraph:P280:begin
\textbf{IV.16.2.}
% mdparagraph:P280:end

% item: EX-IV-16-02
% mdparagraph:P281:begin
Righteousness is more precious than gold.
% mdparagraph:P281:end

\end{minipage}\par\addvspace{4pt}

\par\begin{minipage}{\linewidth}\raggedright
% mdparagraph:P282:begin
\textbf{IV.16.3.}
% mdparagraph:P282:end

% item: EX-IV-16-03
% mdparagraph:P283:begin
The bath is hot.
% mdparagraph:P283:end

\end{minipage}\par\addvspace{4pt}

\par\begin{minipage}{\linewidth}\raggedright
% mdparagraph:P284:begin
\textbf{IV.16.4.}
% mdparagraph:P284:end

% item: EX-IV-16-04
% mdparagraph:P285:begin
The baths are hot.
% mdparagraph:P285:end

\end{minipage}\par\addvspace{4pt}

\par\begin{minipage}{\linewidth}\raggedright
% mdparagraph:P286:begin
\textbf{IV.16.5.}
% mdparagraph:P286:end

% item: EX-IV-16-05
% mdparagraph:P287:begin
The animals are living.
% mdparagraph:P287:end

\end{minipage}\par\addvspace{4pt}

\par\begin{minipage}{\linewidth}\raggedright
% mdparagraph:P288:begin
\textbf{IV.16.6.}
% mdparagraph:P288:end

% item: EX-IV-16-06
% mdparagraph:P289:begin
Some pearls are black.
% mdparagraph:P289:end

\end{minipage}\par\addvspace{4pt}

\par\begin{minipage}{\linewidth}\raggedright
% mdparagraph:P290:begin
\textbf{IV.16.7.}
% mdparagraph:P290:end

% item: EX-IV-16-07
% mdparagraph:P291:begin
The island is beautiful.
% mdparagraph:P291:end

\end{minipage}\par\addvspace{4pt}

\par\begin{minipage}{\linewidth}\raggedright
% mdparagraph:P292:begin
\textbf{IV.16.8.}
% mdparagraph:P292:end

% item: EX-IV-16-08
% mdparagraph:P293:begin
Dolphins swim through the water.
% mdparagraph:P293:end

\end{minipage}\par\addvspace{4pt}

\par\begin{minipage}{\linewidth}\raggedright
% mdparagraph:P294:begin
\textbf{IV.16.9.}
% mdparagraph:P294:end

% item: EX-IV-16-09
% mdparagraph:P295:begin
Here are fishes often caught.
% mdparagraph:P295:end

\end{minipage}\par\addvspace{4pt}

\par\begin{minipage}{\linewidth}\raggedright
% mdparagraph:P296:begin
\textbf{IV.16.10.}
% mdparagraph:P296:end

% item: EX-IV-16-10
% mdparagraph:P297:begin
The night is more beautiful than the day.
% mdparagraph:P297:end

\end{minipage}\par\addvspace{4pt}

\par\begin{minipage}{\linewidth}\raggedright
% mdparagraph:P298:begin
\textbf{IV.16.11.}
% mdparagraph:P298:end

% item: EX-IV-16-11
% mdparagraph:P299:begin
The city was distinguished with walls and gates.
% mdparagraph:P299:end

\end{minipage}\par\addvspace{4pt}

\par\begin{minipage}{\linewidth}\raggedright
% mdparagraph:P300:begin
\textbf{IV.16.12.}
% mdparagraph:P300:end

% item: EX-IV-16-12
% mdparagraph:P301:begin
Summer comes, and then winter.
% mdparagraph:P301:end

\end{minipage}\par\addvspace{4pt}

\par\begin{minipage}{\linewidth}\raggedright
% mdparagraph:P302:begin
\textbf{IV.16.13.}
% mdparagraph:P302:end

% item: EX-IV-16-13
% mdparagraph:P303:begin
The sun shines over all.
% mdparagraph:P303:end

\end{minipage}\par\addvspace{4pt}

\par\begin{minipage}{\linewidth}\raggedright
% mdparagraph:P304:begin
\textbf{IV.16.14.}
% mdparagraph:P304:end

% item: EX-IV-16-14
% mdparagraph:P305:begin
The old man came to the house.
% mdparagraph:P305:end

\end{minipage}\par\addvspace{4pt}

\par\begin{minipage}{\linewidth}\raggedright
% mdparagraph:P306:begin
\textbf{IV.16.15.}
% mdparagraph:P306:end

% item: EX-IV-16-15
% mdparagraph:P307:begin
Here are also found lead and iron.
% mdparagraph:P307:end

\end{minipage}\par\addvspace{4pt}

\par\begin{minipage}{\linewidth}\raggedright
% mdparagraph:P308:begin
\textbf{IV.16.16.}
% mdparagraph:P308:end

% item: EX-IV-16-16
% mdparagraph:P309:begin
Put not your hand into the fire.
% mdparagraph:P309:end

\end{minipage}\par\addvspace{4pt}

\reserve{7\baselineskip}
% mdparagraph:P310:begin
\subsection*{IV.17}
% mdparagraph:P310:end

% mdparagraph:P311:begin
The original “shall” and “will” can have future reference. Keep the actual person and number of the subject. For the predicative “thirsty” use the adjective ðurstig. The numeral for twelve is familiar; the higher count in item17 is supplied separately.
% mdparagraph:P311:end

\reserve{5\baselineskip}
\par\begin{minipage}{\linewidth}\raggedright
% mdparagraph:P312:begin
\textbf{IV.17.1.}
% mdparagraph:P312:end

% item: EX-IV-17-01
% mdparagraph:P313:begin
The soldier takes the gem.
% mdparagraph:P313:end

\end{minipage}\par\addvspace{4pt}

\par\begin{minipage}{\linewidth}\raggedright
% mdparagraph:P314:begin
\textbf{IV.17.2.}
% mdparagraph:P314:end

% item: EX-IV-17-02
% mdparagraph:P315:begin
The rain injures the tower.
% mdparagraph:P315:end

\end{minipage}\par\addvspace{4pt}

\par\begin{minipage}{\linewidth}\raggedright
% mdparagraph:P316:begin
\textbf{IV.17.3.}
% mdparagraph:P316:end

% item: EX-IV-17-03
% mdparagraph:P317:begin
The whale swims.
% mdparagraph:P317:end

\end{minipage}\par\addvspace{4pt}

\par\begin{minipage}{\linewidth}\raggedright
% mdparagraph:P318:begin
\textbf{IV.17.4.}
% mdparagraph:P318:end

% item: EX-IV-17-04
% mdparagraph:P319:begin
Fishermen caught seals.
% mdparagraph:P319:end

\end{minipage}\par\addvspace{4pt}

\par\begin{minipage}{\linewidth}\raggedright
% mdparagraph:P320:begin
\textbf{IV.17.5.}
% mdparagraph:P320:end

% teaching: NEW-U11-LOCAL-HELPS
% teaching: NEW-U11-GRAM-u11-negative-copula-local
% help: U11-COOK-HELP-17-05
% mdparagraph:P321:begin
\textbf{Local help.} For “are not” in this passive, use the supplied plural auxiliary \textbf{ne sindon}. This expression is given for use here; no full negative paradigm is required.
% mdparagraph:P321:end

% item: EX-IV-17-05
% mdparagraph:P322:begin
The hungry sheep are not fed.
% mdparagraph:P322:end

\end{minipage}\par\addvspace{4pt}

\par\begin{minipage}{\linewidth}\raggedright
% mdparagraph:P323:begin
\textbf{IV.17.6.}
% mdparagraph:P323:end

% item: EX-IV-17-06
% mdparagraph:P324:begin
The man’s wife is beautiful.
% mdparagraph:P324:end

\end{minipage}\par\addvspace{4pt}

\par\begin{minipage}{\linewidth}\raggedright
% mdparagraph:P325:begin
\textbf{IV.17.7.}
% mdparagraph:P325:end

% item: EX-IV-17-07
% mdparagraph:P326:begin
Some men are black.
% mdparagraph:P326:end

\end{minipage}\par\addvspace{4pt}

\par\begin{minipage}{\linewidth}\raggedright
% mdparagraph:P327:begin
\textbf{IV.17.8.}
% mdparagraph:P327:end

% item: EX-IV-17-08
% mdparagraph:P328:begin
The carpenter will clothe his son.
% mdparagraph:P328:end

\end{minipage}\par\addvspace{4pt}

\par\begin{minipage}{\linewidth}\raggedright
% mdparagraph:P329:begin
\textbf{IV.17.9.}
% mdparagraph:P329:end

% item: EX-IV-17-09
% mdparagraph:P330:begin
Depart from me, for I am (om. a) sinful man, (om. O) Lord.
% mdparagraph:P330:end

\end{minipage}\par\addvspace{4pt}

\par\begin{minipage}{\linewidth}\raggedright
% mdparagraph:P331:begin
\textbf{IV.17.10.}
% mdparagraph:P331:end

% item: EX-IV-17-10
% mdparagraph:P332:begin
They shall drink, but ye shall be thirsty.
% mdparagraph:P332:end

\end{minipage}\par\addvspace{4pt}

\par\begin{minipage}{\linewidth}\raggedright
% mdparagraph:P333:begin
\textbf{IV.17.11.}
% mdparagraph:P333:end

% help: U11-COOK-HELP-17-11
% mdparagraph:P334:begin
\textbf{Local help.} The give-to-drink expression remains supplied: \textbf{sęle mē wæter drincan}. Sęle is the singular command “give”; mē is the dative recipient and wæter the accusative thing given. The plain infinitive drincan, without tō, says what the recipient is to do.
% mdparagraph:P334:end

% item: EX-IV-17-11
% mdparagraph:P335:begin
He said to her, Give me, I beseech thee, water to drink; for I am thirsty.
% mdparagraph:P335:end

\end{minipage}\par\addvspace{4pt}

\par\begin{minipage}{\linewidth}\raggedright
% mdparagraph:P336:begin
\textbf{IV.17.12.}
% mdparagraph:P336:end

% help: U11-COOK-HELP-17-12
% mdparagraph:P337:begin
\textbf{Local help.} Use supplied \textbf{nis}, “is not”, the third singular negative form of the copula.
% mdparagraph:P337:end

% item: EX-IV-17-12
% mdparagraph:P338:begin
He is not (om. a) God of (om. the) dead, but of (om. the) living.
% mdparagraph:P338:end

\end{minipage}\par\addvspace{4pt}

\par\begin{minipage}{\linewidth}\raggedright
% mdparagraph:P339:begin
\textbf{IV.17.13.}
% mdparagraph:P339:end

% item: EX-IV-17-13
% mdparagraph:P340:begin
The sun shines on (ofer) the righteous and the unrighteous.
% mdparagraph:P340:end

\end{minipage}\par\addvspace{4pt}

\par\begin{minipage}{\linewidth}\raggedright
% mdparagraph:P341:begin
\textbf{IV.17.14.}
% mdparagraph:P341:end

% item: EX-IV-17-14
% mdparagraph:P342:begin
Ye are (om. the) light of the world.
% mdparagraph:P342:end

\end{minipage}\par\addvspace{4pt}

\par\begin{minipage}{\linewidth}\raggedright
% mdparagraph:P343:begin
\textbf{IV.17.15.}
% mdparagraph:P343:end

% item: EX-IV-17-15
% mdparagraph:P344:begin
The twelve gates were twelve pearls.
% mdparagraph:P344:end

\end{minipage}\par\addvspace{4pt}

\par\begin{minipage}{\linewidth}\raggedright
% mdparagraph:P345:begin
\textbf{IV.17.16.}
% mdparagraph:P345:end

% item: EX-IV-17-16
% mdparagraph:P346:begin
The street of the city was gold.
% mdparagraph:P346:end

\end{minipage}\par\addvspace{4pt}

\par\begin{minipage}{\linewidth}\raggedright
% mdparagraph:P347:begin
\textbf{IV.17.17.}
% mdparagraph:P347:end

% teaching: NEW-U11-GRAM-NEW-U11-SRC-NUM4000
% help: U11-COOK-HELP-17-17
% mdparagraph:P348:begin
\textbf{Local help.} Use the supplied phrase \textbf{fēower ðūsend manna}, “four thousand men”. Manna is genitive plural after this higher numeral. Cook’s §154.c concerns this genitive construction. The phrase is supplied as a whole; a general system of higher numerals is not required.
% mdparagraph:P348:end

% item: EX-IV-17-17
% mdparagraph:P349:begin
Those that ate were four thousand men (154. c), besides women and children.
% mdparagraph:P349:end

\end{minipage}\par\addvspace{4pt}

\reserve{7\baselineskip}
% mdparagraph:P350:begin
\subsection*{V.18}
% mdparagraph:P350:end

% mdparagraph:P351:begin
Item1 is a blessing wish with the given past-participle target; keep each title and dependent genitive in its proper case. In item8, the question concerns an actual past event; specify the gender of the singular person addressed as “dearest”. The bracketed “approaches” in item6 explains “draws near”, and (164) refers to the dative construction. The (113) beside hȳdan refers to its weak inflection.
% mdparagraph:P351:end

\reserve{5\baselineskip}
\par\begin{minipage}{\linewidth}\raggedright
% mdparagraph:P352:begin
\textbf{V.18.1.}
% mdparagraph:P352:end

% item: EX-V-18-01
% mdparagraph:P353:begin
Blessed (p.p.) be God, the Father of our Lord Jesus Christ, the Father of mercies, and the God of all comfort.
% mdparagraph:P353:end

\end{minipage}\par\addvspace{4pt}

\par\begin{minipage}{\linewidth}\raggedright
% mdparagraph:P354:begin
\textbf{V.18.2.}
% mdparagraph:P354:end

% item: EX-V-18-02
% mdparagraph:P355:begin
The winds were favorable.
% mdparagraph:P355:end

\end{minipage}\par\addvspace{4pt}

\par\begin{minipage}{\linewidth}\raggedright
% mdparagraph:P356:begin
\textbf{V.18.3.}
% mdparagraph:P356:end

% item: EX-V-18-03
% mdparagraph:P357:begin
He wishes to watch their journey.
% mdparagraph:P357:end

\end{minipage}\par\addvspace{4pt}

\par\begin{minipage}{\linewidth}\raggedright
% mdparagraph:P358:begin
\textbf{V.18.4.}
% mdparagraph:P358:end

% item: EX-V-18-04
% mdparagraph:P359:begin
The wind and the sea obey him.
% mdparagraph:P359:end

\end{minipage}\par\addvspace{4pt}

\par\begin{minipage}{\linewidth}\raggedright
% mdparagraph:P360:begin
\textbf{V.18.5.}
% mdparagraph:P360:end

% item: EX-V-18-05
% mdparagraph:P361:begin
Not at all in vain fought I with death, as he fought with beasts.
% mdparagraph:P361:end

\end{minipage}\par\addvspace{4pt}

\par\begin{minipage}{\linewidth}\raggedright
% mdparagraph:P362:begin
\textbf{V.18.6.}
% mdparagraph:P362:end

% item: EX-V-18-06
% mdparagraph:P363:begin
The time draws near [approaches, 164] (om. the) birth of Christ.
% mdparagraph:P363:end

\end{minipage}\par\addvspace{4pt}

\par\begin{minipage}{\linewidth}\raggedright
% mdparagraph:P364:begin
\textbf{V.18.7.}
% mdparagraph:P364:end

% teaching: NEW-U11-GRAM-u11-birth-passive-supplied
% help: U11-COOK-HELP-18-07
% mdparagraph:P365:begin
\textbf{Local help.} Cook supplies \textbf{ācenned}, “born”, a past participle, with plain e. Use it with the past passive auxiliary for “was born”; the other forms of this verb are not required.
% mdparagraph:P365:end

% item: EX-V-18-07
% mdparagraph:P366:begin
It is the day when (ðe) he was born (ācenned).
% mdparagraph:P366:end

\end{minipage}\par\addvspace{4pt}

\par\begin{minipage}{\linewidth}\raggedright
% mdparagraph:P367:begin
\textbf{V.18.8.}
% mdparagraph:P367:end

% item: EX-V-18-08
% mdparagraph:P368:begin
O, wast thou with me, dearest, then, while I rose up against my doom (dōm)?
% mdparagraph:P368:end

\end{minipage}\par\addvspace{4pt}

\par\begin{minipage}{\linewidth}\raggedright
% mdparagraph:P369:begin
\textbf{V.18.9.}
% mdparagraph:P369:end

% item: EX-V-18-09
% mdparagraph:P370:begin
The moon is hid (hȳdan, 113), the night is still.
% mdparagraph:P370:end

\end{minipage}\par\addvspace{4pt}

\par\begin{minipage}{\linewidth}\raggedright
% mdparagraph:P371:begin
\textbf{V.18.10.}
% mdparagraph:P371:end

% item: EX-V-18-10
% mdparagraph:P372:begin
Thou bringst the fisher to his wife.
% mdparagraph:P372:end

\end{minipage}\par\addvspace{4pt}

\par\begin{minipage}{\linewidth}\raggedright
% mdparagraph:P373:begin
\textbf{V.18.11.}
% mdparagraph:P373:end

% item: EX-V-18-11
% mdparagraph:P374:begin
Calm is the morn(ing).
% mdparagraph:P374:end

\end{minipage}\par\addvspace{4pt}

\par\begin{minipage}{\linewidth}\raggedright
% mdparagraph:P375:begin
\textbf{V.18.12.}
% mdparagraph:P375:end

% item: EX-V-18-12
% mdparagraph:P376:begin
I know that this was life.
% mdparagraph:P376:end

\end{minipage}\par\addvspace{4pt}

\par\begin{minipage}{\linewidth}\raggedright
% mdparagraph:P377:begin
\textbf{V.18.13.}
% mdparagraph:P377:end

% item: EX-V-18-13
% mdparagraph:P378:begin
What do ye more than they?
% mdparagraph:P378:end

\end{minipage}\par\addvspace{4pt}

\reserve{7\baselineskip}
% mdparagraph:P379:begin
\section*{12. The fisherman and his son}
% mdparagraph:P379:end

% mdparagraph:P380:begin
This connected account is new teaching material. Read it as a sequence: arrival, assistance, speech and the action that follows. Resolve pronouns from both their forms and the context.
% mdparagraph:P380:end

% composition: NEW-U11-READING-01
% mdparagraph:P381:begin
\par\begin{minipage}{\linewidth}\raggedright
\readingtext{Se ealda fiscere cōm tō ðǣm hūse, ðæt hē his suna hulpe. Ðā hē cōm, ðā wæs se sunu hungrig. Se fiscere sealde him męte. Mid-ðȳ se sunu ǣt, se fiscere bær wæter. Ðēah se fiscere wǣre eald, hē healp his suna. Se sunu cwæð: “Wes hēr, oð-ðæt ic cume.” Hē gewāt and cōm and sealde his fæder wæter.}
\end{minipage}\par
% mdparagraph:P381:end

\reserve{5\baselineskip}
\par\begin{minipage}{\linewidth}\raggedright
% mdparagraph:P382:begin
\textbf{R01.}
% mdparagraph:P382:end

% item: NEW-U11-R01
% mdparagraph:P383:begin
Translate the whole new account The fisherman and his son.
% mdparagraph:P383:end

\end{minipage}\par\addvspace{4pt}

\par\begin{minipage}{\linewidth}\raggedright
% mdparagraph:P384:begin
\textbf{R02.}
% mdparagraph:P384:end

% item: NEW-U11-R02
% mdparagraph:P385:begin
In the new account, identify the purpose clause, the while clause and the concession. State what each contributes to the account and parse its finite verb.
% mdparagraph:P385:end

\end{minipage}\par\addvspace{4pt}

\par\begin{minipage}{\linewidth}\raggedright
% mdparagraph:P386:begin
\textbf{R03.}
% mdparagraph:P386:end

% item: NEW-U11-R03
% mdparagraph:P387:begin
In the new account, who gives the first food, who receives it, and who gives the final water? Explain the relevant pronouns and genitive possessors.
% mdparagraph:P387:end

\end{minipage}\par\addvspace{4pt}

\par\begin{minipage}{\linewidth}\raggedright
% mdparagraph:P388:begin
\textbf{R04.}
% mdparagraph:P388:end

% item: NEW-U11-R04
% mdparagraph:P389:begin
In the son’s direct speech, who is ic, when is the coming expected, and what mood is cume? Identify the later action that fulfills the expectation.
% mdparagraph:P389:end

\end{minipage}\par\addvspace{4pt}

\reserve{7\baselineskip}
% mdparagraph:P390:begin
\section*{13. Bede’s description of Britain: metals and jet}
% mdparagraph:P390:end

% mdparagraph:P391:begin
This is an exact extract from Cook, \emph{A First Book in Old English}, printed p.139 /\allowbreak{} PDF p.155, editorial lines1–4, ending at onwęg. Cook’s attribution is “(Eccl. Hist., Bk. I., Chap. I.)”. The passage has 35 words; the dash is punctuation. The line count is editorial, anchored by the printed marginal line5. Original notes are preserved in §15.
% mdparagraph:P391:end

% mdparagraph:P392:begin
\textbf{Context.} The subject is Britain. Bede’s preceding description has discussed the island’s plants, animals, shellfish and springs.
% mdparagraph:P392:end

\par\begin{minipage}{\linewidth}\raggedright
% source: FB-IV-BRITAIN-METALS-01
% mdparagraph:P393:begin
\readingtext{Swilce hit is ēac berende on węcga ōrum — āres and īsernes, lēades and seolfres. Hēr bið ēac gemētt gagates; se stān bið blæc gimm; gif man hine on fȳr dēð, ðonne flēoð ðǣr næddran onwęg.}
% mdparagraph:P393:end

\end{minipage}\par
\reserve{7\baselineskip}
% mdparagraph:P394:begin
\subsection*{Supplied help}
% mdparagraph:P394:end

% mdparagraph:P395:begin
The following words and constructions may be used with their help throughout the reading tasks. The whole expressions are supported here; they need not be reconstructed from earlier headwords.
% mdparagraph:P395:end

% teaching: NEW-U11-GRAM-NEW-U11-SRC-B01
% help: U11-SRC-HELP-B01
% mdparagraph:P396:begin
\textbf{1. Swilce hit ... seolfres.} The preceding context concerns Britain, the island: hit refers to it. Swilce ... ēac, also. Berende on, rich in /\allowbreak{} productive of; węcga is genitive plural “of metals,” and ōrum dative plural “veins/\allowbreak{}ores” after on. The following genitives name the metals. Āres is Cook’s source spelling for copper here, beside the earlier selected ǣr.
% mdparagraph:P396:end

% teaching: NEW-U11-GRAM-NEW-U11-SRC-B02
% help: U11-SRC-HELP-B02
% mdparagraph:P397:begin
\textbf{2. Hēr bið ... blæc gimm.} Gemētt, found, is the exact past participle of prefixed gemētan; bið gemētt, is found. Gagates is jet, the stone’s Latin loan-name; stān, stone, masculine nominative singular. The second bið links stone with black gem.
% mdparagraph:P397:end

% teaching: NEW-U11-GRAM-NEW-U11-SRC-B03
% help: U11-SRC-HELP-B03
% mdparagraph:P398:begin
\textbf{3. gif man hine on fȳr dēð, ðonne.} Man is the impersonal singular subject “one”; dēð is third singular present indicative. Hine refers to the masculine stone. The newly taught placement frame dōn + object + on with accusative gives “puts it on/\allowbreak{}in a fire”; ðonne correlates with the gif condition.
% mdparagraph:P398:end

% teaching: NEW-U11-GRAM-NEW-U11-SRC-B04
% help: U11-SRC-HELP-B04
% mdparagraph:P399:begin
\textbf{4. flēoð ðǣr næddran onwęg.} Flēoð is third plural present indicative of flēon, flee; it is not the verb fly. Ðǣr, there; næddran, snakes/\allowbreak{}adders, nominative plural of weak feminine næddre; onwęg, away. Both d letters in næddran belong to the printed text.
% mdparagraph:P399:end

\reserve{5\baselineskip}
\par\begin{minipage}{\linewidth}\raggedright
% mdparagraph:P400:begin
\textbf{B01.}
% mdparagraph:P400:end

% item: NEW-U11-B01
% mdparagraph:P401:begin
With the printed helps, translate the complete Bede metals-and-jet passage.
% mdparagraph:P401:end

\end{minipage}\par\addvspace{4pt}

\par\begin{minipage}{\linewidth}\raggedright
% mdparagraph:P402:begin
\textbf{B02.}
% mdparagraph:P402:end

% item: NEW-U11-B02
% mdparagraph:P403:begin
In the Bede passage, identify the condition introduced by gif and its main clause. Parse the two finite verbs using the supplied forms, and explain the role of ðonne.
% mdparagraph:P403:end

\end{minipage}\par\addvspace{4pt}

\par\begin{minipage}{\linewidth}\raggedright
% mdparagraph:P404:begin
\textbf{B03.}
% mdparagraph:P404:end

% item: NEW-U11-B03
% mdparagraph:P405:begin
What does the first hit refer to, and what does hine in the condition refer to? Give the grammatical information that distinguishes them, using the introduction and helps.
% mdparagraph:P405:end

\end{minipage}\par\addvspace{4pt}

\reserve{7\baselineskip}
% mdparagraph:P406:begin
\section*{14. Æthelwald calms the sea}
% mdparagraph:P406:end

% mdparagraph:P407:begin
This exact connected episode is from Cook’s \emph{First Book}, printed pp.142–143 /\allowbreak{} PDF pp.158–159: editorial line8 on p.142 through line14 on p.143. It begins after standan. in the preceding account and ends with heofonlice forgiefen wæs. The 195 words retain the original two paragraphs. Cook’s original notes, including the cross-reference to an earlier note, follow in §15.
% mdparagraph:P407:end

% mdparagraph:P408:begin
\textbf{Earlier events.} Three visitors are returning by boat from the hermit Æthelwald on Farne. A sudden storm has stopped their progress. After struggling in vain, they look back toward the island. The narrator now tells how they were rescued.
% mdparagraph:P408:end

\par\begin{minipage}{\linewidth}\raggedright
% source: FB-V-SEA-RESCUE-01
% mdparagraph:P409:begin
\readingtext{Ðā wæs æfter langum fæce ðæt wē ūre gesihð feorr ūpāhōfon, ðā gesāwon wē in Farne, ðǣm īglande, Gode ðone lēofan fæder Æðelwald of his dīegelnessum ūtgangende, ðæt hē wolde ūrne sīðfæt scēawian, and gesēon hwæt ūs gelumpe, for-ðon hē gehīerde ðæt gebrec ðǣra storma and ðæs weallendes sǣs.}
% mdparagraph:P409:end

% mdparagraph:P410:begin
\readingtext{Mid-ðȳ hē ðā ūs ēac scēawode, and geseah in gewinne and in ōrmōdnesse gesętte bēon, ðā biegde hē his cnēowu tō Fæder ūres Dryhtnes Hǣlendes Crīstes, and wæs gebiddende for ūre hǣle and for ūrum līfe. And mid-ðȳ hē ðā ðæt gebed gefylde, hē ðā samod ætgædere ge ðone āðundnan sǣ gesmylte ge ðone storm gestilde, tō ðon ðætte ðurh eall sēo rēðnes ðæs stormes wæs blinnende, and gesyndge windas ðurh ðone smyltestan sǣ ūs æt lande gebrōhton. Mid-ðȳ wē ðā ūp cōmon tō lande, and ūre scip ēac swilce fram ðǣm ȳðum ūp ābǣron, ðā sōna se ilca storm ęft hwearf and cōm, se-ðe for ūrum intingan medmicel fæc gestilde, and ealne ðone dæg swīðe micel and strang wæs, ðætte męnn sweotollice ongietan meahton ðætte se medmicla fierst ðǣre stilnesse, ðe ðǣr becōm, tō bēnum ðæs Godes weres for intingan ūre hǣle heofonlice forgiefen wæs.}
% mdparagraph:P410:end

\end{minipage}\par
\reserve{7\baselineskip}
% mdparagraph:P411:begin
\subsection*{Supplied help}
% mdparagraph:P411:end

% mdparagraph:P412:begin
Follow the shifts from what the travellers see to Æthelwald’s actions and back to the travellers. The complete expressions below are supplied. The sea remains masculine in these phrases, as in the earlier paradigm.
% mdparagraph:P412:end

% teaching: NEW-U11-GRAM-NEW-U11-SRC-S01
% help: U11-SRC-HELP-S01
% mdparagraph:P413:begin
\textbf{1. Ðā wæs æfter langum fæce ðæt.} Read the whole opening as “then, after a long time, it happened that ...”. Æfter governs the dative; langum fæce means a long interval. The opening ðæt belongs to this supplied time-introduction frame.
% mdparagraph:P413:end

% teaching: NEW-U11-GRAM-NEW-U11-SRC-S02
% help: U11-SRC-HELP-S02
% mdparagraph:P414:begin
\textbf{2. wē ūre gesihð feorr ūpāhōfon.} Gesihð is feminine accusative singular, sight/\allowbreak{}gaze; ūre agrees with it. Feorr, far off. Ūpāhōfon, we lifted up, first plural preterite of ūpāhebban. The complete phrase describes looking up and seeing at a distance.
% mdparagraph:P414:end

% teaching: NEW-U11-GRAM-NEW-U11-SRC-S03
% help: U11-SRC-HELP-S03
% mdparagraph:P415:begin
\textbf{3. in Farne ... of his dīegelnessum ūtgangende.} In Farne, on Farne island, with ðǣm īglande as dative apposition. Gode ðone lēofan fæder means the father dear to God: the adjective’s complement is dative. Æðelwald is in apposition to the accusative fæder. Of + dative means out of/\allowbreak{}from; dīegelnessum is the source plural for retreat/\allowbreak{}seclusion. Ūtgangende is the printed participle, coming out; Cook notes that a present participle is used for a Latin past participle here, so “having come out” is also possible.
% mdparagraph:P415:end

% teaching: NEW-U11-GRAM-NEW-U11-SRC-S04
% help: U11-SRC-HELP-S04
% mdparagraph:P416:begin
\textbf{4. ðæt hē wolde ... hwæt ūs gelumpe.} The ðæt clause gives his purpose. Wolde with the infinitives scēawian and gesēon is read here “that he might watch and see” or “in order to watch and see”; do not make private desire the only force. Hwæt ūs gelumpe is a dependent question, what might happen to us. Gelumpe is third singular preterite subjunctive, and ūs is dative.
% mdparagraph:P416:end

% teaching: NEW-U11-GRAM-NEW-U11-SRC-S05
% help: U11-SRC-HELP-S05
% mdparagraph:P417:begin
\textbf{5. for-ðon hē gehīerde ... weallendes sǣs.} For-ðon is the shorter source spelling of the causal conjunction. Gehīerde, heard, is a prefixed past form. Gebrec is the accusative noise/\allowbreak{}crash; ðǣra storma is genitive plural of storms, and ðæs weallendes sǣs genitive singular of the surging sea. Weallendes is an adjectival participle, surging/\allowbreak{}boiling. Ðǣra preserves the source genitive-plural spelling.
% mdparagraph:P417:end

% teaching: NEW-U11-GRAM-NEW-U11-SRC-S06
% help: U11-SRC-HELP-S06
% mdparagraph:P418:begin
\textbf{6. Mid-ðȳ ... in gewinne and in ōrmōdnesse gesętte bēon.} The temporal clause contains both scēawode and geseah. With geseah, understand the earlier ūs as the people seen: he saw us to be situated in distress and despair. In governs the two datives gewinne and ōrmōdnesse. Gesętte is plural past participle here, not the identically spelled finite past; bēon is infinitive. This entire perception complement is supplied.
% mdparagraph:P418:end

% teaching: NEW-U11-GRAM-NEW-U11-SRC-S07
% help: U11-SRC-HELP-S07
% mdparagraph:P419:begin
\textbf{7. ðā biegde hē his cnēowu tō Fæder ... Crīstes.} Biegde, bent, is the exact third singular preterite spelling; cnēowu, knees, accusative plural. Tō Fæder marks the direction of the prayer. Ūres Dryhtnes Hǣlendes Crīstes is a genitive group: of our Lord Jesus Christ, with the proper title Hǣlend and Crīst in apposition.
% mdparagraph:P419:end

% teaching: NEW-U11-GRAM-NEW-U11-SRC-S08
% help: U11-SRC-HELP-S08
% mdparagraph:P420:begin
\textbf{8. wæs gebiddende for ūre hǣle and for ūrum līfe.} Cook143 note1 refers to142 note7: this was + participle wording is read as a simple past here, “prayed.” Do not infer a universal English progressive rule. For ūre hǣle, for our safety/\allowbreak{}rescue, uses accusative; for ūrum līfe, for our life, uses dative. Keep both exact constructions supplied.
% mdparagraph:P420:end

% teaching: NEW-U11-GRAM-NEW-U11-SRC-S09
% help: U11-SRC-HELP-S09
% mdparagraph:P421:begin
\textbf{9. mid-ðȳ ... ðæt gebed gefylde ... samod ætgædere ge.} Gebed is prayer, accusative singular; gefylde is the already taught finished/\allowbreak{}accomplished. Samod ætgædere means together/\allowbreak{}at once. Paired ge ... ge coordinates both ... and; the first ge introduces the sea-calming and the second the storm-stilling. The narrated hē remains Æthelwald, the person who has just completed the prayer.
% mdparagraph:P421:end

% teaching: NEW-U11-GRAM-NEW-U11-SRC-S10
% help: U11-SRC-HELP-S10
% mdparagraph:P422:begin
\textbf{10. ðone āðundnan sǣ gesmylte ge ðone storm gestilde.} Āðundnan, swollen, is a weak participial adjective agreeing with masculine accusative sǣ. Gesmylte, calmed, is third singular preterite; storm is accusative singular. Gestilde retains the taught transitive quell/\allowbreak{}settle sense here.
% mdparagraph:P422:end

% teaching: NEW-U11-GRAM-NEW-U11-SRC-S11
% help: U11-SRC-HELP-S11
% mdparagraph:P423:begin
\textbf{11. tō ðon ðætte ðurh eall sēo rēðnes ... wæs blinnende.} Tō ðon ðætte introduces the result, so that. Ðurh eall is the exact adverbial expression completely/\allowbreak{}entirely. Rēðnes, fury, is the feminine subject; ðæs stormes is genitive. Wæs blinnende is the supplied periphrasis “ceased” here, with present participle blinnende; no new productive progressive tense is introduced.
% mdparagraph:P423:end

% teaching: NEW-U11-GRAM-NEW-U11-SRC-S12
% help: U11-SRC-HELP-S12
% mdparagraph:P424:begin
\textbf{12. gesyndge windas ... æt lande gebrōhton.} Gesyndge is the exact syncopated source form beside taught gesyndige, favorable, nominative plural. Smyltestan is the weak masculine accusative superlative calmest/\allowbreak{}very calm. Ðurh governs ðone ... sǣ; ūs is the accusative people brought. Æt lande means to/\allowbreak{}at land in this exact arrival expression. Gebrōhton is the prefixed third plural preterite brought.
% mdparagraph:P424:end

% teaching: NEW-U11-GRAM-NEW-U11-SRC-S13
% help: U11-SRC-HELP-S13
% mdparagraph:P425:begin
\textbf{13. Mid-ðȳ wē ... fram ðǣm ȳðum ūp ābǣron.} The temporal clause has the two completed actions: we came up onto land and carried our ship up from the waves. Lande is the locally supplied dative of land; ȳðum is dative plural waves after fram. Ābǣron is first plural preterite of prefixed carry. Ūre scip is the accusative object, not a new subject.
% mdparagraph:P425:end

% teaching: NEW-U11-GRAM-NEW-U11-SRC-S14
% help: U11-SRC-HELP-S14
% mdparagraph:P426:begin
\textbf{14. ðā sōna ... se-ðe for ūrum intingan medmicel fæc gestilde.} Sōna, immediately; ęft, again; hwearf, returned. Se-ðe is the exact relative complex referring to masculine storm. For ūrum intingan, for our sake; intingan is dative singular. Medmicel fæc is accusative duration, a short interval. Gestilde is intransitive ceased/\allowbreak{}became still in this relative clause; the earlier transitive sense alone is insufficient. English had ceased may clarify the sequence.
% mdparagraph:P426:end

% teaching: NEW-U11-GRAM-NEW-U11-SRC-S15
% help: U11-SRC-HELP-S15
% mdparagraph:P427:begin
\textbf{15. and ealne ðone dæg ... ðætte męnn sweotollice ongietan meahton ðætte.} Ealne ðone dæg is accusative duration, all that day. Strang, strong, agrees with storm. The first ðætte introduces the result: people could clearly understand. Sweotollice, clearly; ongietan, understand, is the infinitive governed by meahton. The second ðætte introduces what they understand, a content clause.
% mdparagraph:P427:end

% teaching: NEW-U11-GRAM-NEW-U11-SRC-S16
% help: U11-SRC-HELP-S16
% mdparagraph:P428:begin
\textbf{16. se medmicla fierst ðǣre stilnesse, ðe ðǣr becōm.} Fierst is masculine nominative singular, interval/\allowbreak{}period; medmicla is its weak adjective, short. Ðǣre stilnesse, of the calm, is genitive singular. In the relative clause, ðǣr means there and becōm means occurred/\allowbreak{}came about in this context. Relative ðe is the subject of becōm and refers to the interval; it is distinct from the following adverb ðǣr.
% mdparagraph:P428:end

% teaching: NEW-U11-GRAM-NEW-U11-SRC-S17
% help: U11-SRC-HELP-S17
% mdparagraph:P429:begin
\textbf{17. tō bēnum ðæs Godes weres ... heofonlice forgiefen wæs.} Tō bēnum, at/\allowbreak{}in response to the prayers, uses dative plural. Ðæs Godes weres means of the man of God; weres is genitive singular, with Godes depending on it. For intingan ūre hǣle means for the sake of our rescue: hǣle is genitive singular; exact ūre is supplied in this possessive expression. Heofonlice means from heaven; forgiefen means given/\allowbreak{}granted here, expressly not forgiven. The final passive refers to the brief period of calm. Here ūre is the first-person plural genitive used attributively within “our rescue”; this supplied phrase does not establish a new feminine possessive-adjective variant.
% mdparagraph:P429:end

\reserve{5\baselineskip}
\par\begin{minipage}{\linewidth}\raggedright
% mdparagraph:P430:begin
\textbf{S01.}
% mdparagraph:P430:end

% item: NEW-U11-S01
% mdparagraph:P431:begin
With the printed helps, translate the complete rescue-and-resolution passage from Æthelwald Calms the Sea.
% mdparagraph:P431:end

\end{minipage}\par\addvspace{4pt}

\par\begin{minipage}{\linewidth}\raggedright
% mdparagraph:P432:begin
\textbf{S02.}
% mdparagraph:P432:end

% item: NEW-U11-S02
% mdparagraph:P433:begin
Put these events in the order given by the passage: the storm returns; the sailors see Æthelwald; the sailors reach land; Æthelwald prays; the sea becomes calm. State what the sailors have done before the storm returns.
% mdparagraph:P433:end

\end{minipage}\par\addvspace{4pt}

\par\begin{minipage}{\linewidth}\raggedright
% mdparagraph:P434:begin
\textbf{S03.}
% mdparagraph:P434:end

% item: NEW-U11-S03
% mdparagraph:P435:begin
Find the clause that gives Æthelwald’s purpose in coming out and the clause that gives his reason. Explain their different functions; use the local help for wolde and its infinitives.
% mdparagraph:P435:end

\end{minipage}\par\addvspace{4pt}

\par\begin{minipage}{\linewidth}\raggedright
% mdparagraph:P436:begin
\textbf{S04.}
% mdparagraph:P436:end

% item: NEW-U11-S04
% mdparagraph:P437:begin
Find the result introduced by tō ðon ðætte. What happened to the storm, and what brought the sailors to land? Use the help to explain wæs blinnende.
% mdparagraph:P437:end

\end{minipage}\par\addvspace{4pt}

\par\begin{minipage}{\linewidth}\raggedright
% mdparagraph:P438:begin
\textbf{S05.}
% mdparagraph:P438:end

% item: NEW-U11-S05
% mdparagraph:P439:begin
Identify the two mid-ðȳ clauses that introduce Æthelwald’s sight of the sailors and the sailors’ coming ashore. Where does the main action begin in each case?
% mdparagraph:P439:end

\end{minipage}\par\addvspace{4pt}

\par\begin{minipage}{\linewidth}\raggedright
% mdparagraph:P440:begin
\textbf{S06.}
% mdparagraph:P440:end

% item: NEW-U11-S06
% mdparagraph:P441:begin
Explain how wæs gebiddende is to be understood here, using Cook’s note and the help. Why does the present participle here not require an English progressive translation?
% mdparagraph:P441:end

\end{minipage}\par\addvspace{4pt}

\par\begin{minipage}{\linewidth}\raggedright
% mdparagraph:P442:begin
\textbf{S07.}
% mdparagraph:P442:end

% item: NEW-U11-S07
% mdparagraph:P443:begin
At the end, what do people understand about the short period of calm? Say whose prayers and whose rescue are named, and explain the final forgiefen using the supplied note.
% mdparagraph:P443:end

\end{minipage}\par\addvspace{4pt}

\par\begin{minipage}{\linewidth}\raggedright
% mdparagraph:P444:begin
\textbf{S08.}
% mdparagraph:P444:end

% item: NEW-U11-S08
% mdparagraph:P445:begin
Does the passage say that the storm ceased permanently? Explain the return of the storm and how it supports the narrator’s final explanation.
% mdparagraph:P445:end

\end{minipage}\par\addvspace{4pt}

\reserve{8\baselineskip}
% mdparagraph:P446:begin
\section*{15. Cook’s original notes and source information}
% mdparagraph:P446:end

\begingroup\fontsize{10}{12.8}\selectfont
\setlength{\parskip}{1.5pt}
% mdparagraph:P447:begin
These are Cook’s endnotes, kept separate from the Old English and the new learner help. Numbers belong to their printed pages; they restart on each page. His historical and etymological questions remain source apparatus and are not additional exercises. The anchors identify the words bearing the original cues.
% mdparagraph:P447:end

\reserve{7\baselineskip}
% mdparagraph:P448:begin
\subsection*{Printed page 139}
% mdparagraph:P448:end

% source-note: U11-SRC-P139-N01
% mdparagraph:P449:begin
\textbf{1. berende.} Cf. the construction of this word with that above, p. 138, l. 1.
% mdparagraph:P449:end

% source-note: U11-SRC-P139-N02
% mdparagraph:P450:begin
\textbf{2. man.} See 89. e.
% mdparagraph:P450:end

% source-note: U11-SRC-P139-N03
% mdparagraph:P451:begin
\textbf{3. dēð.} See 140.
% mdparagraph:P451:end

% source-note: U11-SRC-P139-N04
% mdparagraph:P452:begin
\textbf{4. næddran.} How did næddre become adder? Cf. OE. nafogār, Mod. Eng. auger. See Skeat, Prin., p. 216.
% mdparagraph:P452:end

% source-note: U11-SRC-P139-N05
% mdparagraph:P453:begin
\textbf{5. onwęg.} There is a parallel form, āwęg, already in OE. The ā- is a contraction of on. Mention other Mod. Eng. words in which the a- represents on.
% mdparagraph:P453:end

\reserve{7\baselineskip}
% mdparagraph:P454:begin
\subsection*{Printed page 142}
% mdparagraph:P454:end

% source-note: U11-SRC-P142-N14
% mdparagraph:P455:begin
\textbf{14. feorr.} Translate, from a distance.
% mdparagraph:P455:end

% source-note: U11-SRC-P142-N15
% mdparagraph:P456:begin
\textbf{15. Gode.} Governed by lēofan; = amantissimum Deo. See 165.
% mdparagraph:P456:end

% source-note: U11-SRC-P142-N16
% mdparagraph:P457:begin
\textbf{16. dīegelnessum.} Lat. latibulis.
% mdparagraph:P457:end

% source-note: U11-SRC-P142-N17
% mdparagraph:P458:begin
\textbf{17. ūtgangende.} Translates the Latin past part., egressum.
% mdparagraph:P458:end

% source-note: U11-SRC-P142-N18
% mdparagraph:P459:begin
\textbf{18. ðæt hē wolde.} Translate, that he might, in order to, to. The Latin has the infinitive.
% mdparagraph:P459:end

% source-note: U11-SRC-P142-N19
% mdparagraph:P460:begin
\textbf{19. gelumpe.} See 194. b.
% mdparagraph:P460:end

% source-note: U11-SRC-P142-N20
% mdparagraph:P461:begin
\textbf{20. sǣs.} Lat. fragore procellarum ac ferventis oceani.
% mdparagraph:P461:end

% source-note: U11-SRC-P142-N21
% mdparagraph:P462:begin
\textbf{21. ōrmōdnesse.} Lat. desperatione.
% mdparagraph:P462:end

% source-note: U11-SRC-P142-N22
% mdparagraph:P463:begin
\textbf{22. bēon.} Supply ūs as subject acc. (169).
% mdparagraph:P463:end

\reserve{7\baselineskip}
% mdparagraph:P464:begin
\subsection*{Printed page 143}
% mdparagraph:P464:end

% source-note: U11-SRC-P143-N01
% mdparagraph:P465:begin
\textbf{1. wæs gebiddende; wæs blinnende.} See above, p. 142, n. 7.
% mdparagraph:P465:end

% source-note: U11-SRC-P143-N02
% mdparagraph:P466:begin
\textbf{2. gefylde.} Lat. compleret.
% mdparagraph:P466:end

% source-note: U11-SRC-P143-N03
% mdparagraph:P467:begin
\textbf{3. āðundnan.} Lat. tumida.
% mdparagraph:P467:end

% source-note: U11-SRC-P143-N04
% mdparagraph:P468:begin
\textbf{4. tō ðon ðætte.} Lat. adeo ut, nearly = so that.
% mdparagraph:P468:end

% source-note: U11-SRC-P143-N05
% mdparagraph:P469:begin
\textbf{5. ðurh eall.} Lat. per omnia, = entirely.
% mdparagraph:P469:end

% source-note: U11-SRC-P143-N06
% mdparagraph:P470:begin
\textbf{6. gesyndge.} Lat. secunda. What letter (sound) has been lost from the OE. form?
% mdparagraph:P470:end

% source-note: U11-SRC-P143-N07
% mdparagraph:P471:begin
\textbf{7. ēac swilce.} Lat. quoque.
% mdparagraph:P471:end

% source-note: U11-SRC-P143-N08
% mdparagraph:P472:begin
\textbf{8. cōm.} Supply ęft. Latin has only one verb, rediit.
% mdparagraph:P472:end

% source-note: U11-SRC-P143-N09
% mdparagraph:P473:begin
\textbf{9. for ūrum intingan.} Lat. nostri gratia.
% mdparagraph:P473:end

% source-note: U11-SRC-P143-N10
% mdparagraph:P474:begin
\textbf{10. fæc; dæg.} See 170.
% mdparagraph:P474:end

% source-note: U11-SRC-P143-N11
% mdparagraph:P475:begin
\textbf{11. ðone.} Lat. illo. Translate, that.
% mdparagraph:P475:end

% source-note: U11-SRC-P143-N12
% mdparagraph:P476:begin
\textbf{12. ðætte.} From ðæt-ðe (34).
% mdparagraph:P476:end

% source-note: U11-SRC-P143-N13
% mdparagraph:P477:begin
\textbf{13. tō bēnum.} Lat. ad preces. Cf. the Mod. Eng. phrase, ‘bootless bene.’
% mdparagraph:P477:end

% source-note: U11-SRC-P143-N14
% mdparagraph:P478:begin
\textbf{14. Godes.} Dependent on weres.
% mdparagraph:P478:end

% source-note: U11-SRC-P143-N15
% mdparagraph:P479:begin
\textbf{15. weres.} Dependent on bēnum.
% mdparagraph:P479:end

% source-note: U11-SRC-P143-N16
% mdparagraph:P480:begin
\textbf{16. hǣle.} Lat. evasionis.
% mdparagraph:P480:end

% source-note: U11-SRC-P143-N17
% mdparagraph:P481:begin
\textbf{17. heofonlice.} Lat. cælitus, = from heaven.
% mdparagraph:P481:end

% source-note: U11-SRC-P143-N18
% mdparagraph:P482:begin
\textbf{18. forgiefen.} Not forgiven, but given.
% mdparagraph:P482:end

\reserve{7\baselineskip}
% mdparagraph:P483:begin
\subsection*{Referenced note: printed page 142, note 7}
% mdparagraph:P483:end

% mdparagraph:P484:begin
Cook p.143 note1 explicitly refers to this earlier explanation. The preceding sentence to which it was first attached is not part of the assigned extract.
% mdparagraph:P484:end

% source-note: U11-SRC-P142-N07
% mdparagraph:P485:begin
Lat. egressi eramus. The pres. part. with the verb is sometimes used in OE. to denote the simple past, as here, and not the progressive.
% mdparagraph:P485:end

\reserve{7\baselineskip}
% mdparagraph:P486:begin
\subsection*{Editions, spelling and grammatical references}
% mdparagraph:P486:end

% mdparagraph:P487:begin
Cook’s \emph{Exercises in Old English} (1895), printed pp.8–9 /\allowbreak{} one-based PDF pp.16–17, supplies all sixteen IV.16 items, all seventeen IV.17 items and all thirteen V.18 items. The original references (64), (154. c), (164) and (113), the omission directions, and the hints dōm and ācenned are retained. The Exercises page offset is +8; for Cook’s \emph{A First Book in Old English}, second edition (1895), the offset is +16. Printed pages and editorial line counts are distinct from PDF pages.
% mdparagraph:P487:end

% mdparagraph:P488:begin
In the exact readings, printed line wraps and line-end divisions are joined and note cues are placed in the endnotes; lexical hyphens, letters, quantities and paragraph divisions remain. Retain āres, næddran, ðǣra, biegde, gesyndge and se-ðe as printed. The first passage describes a belief about jet; translating the statement does not turn it into a modern scientific claim.
% mdparagraph:P488:end

% mdparagraph:P489:begin
The noun patterns and their exceptions follow Cook §§43–54, pp.28–38; adjective patterns §§55–64, pp.38–43. The reduction in comparative genitive plurals continues the earlier rule; compare Wright, \emph{Old English Grammar} (1914), §259, p.130. Selected vocabulary entries establish the genders, stems and meanings. The proper-title construction follows Cook §151, p.88. Personal adjectives follow §§55 and181. The participle bases and indicative systems follow the relevant strong and weak paradigms; complete tables here combine those patterns with the individual lexical stems.
% mdparagraph:P489:end

% mdparagraph:P490:begin
For the clause relations compare Cook §§194–196, pp.103–105, with Fulk, \emph{An Introductory Grammar of Old English} (2020), §45. Constructions with infinitives and perception follow Cook §§169 and198–200; the exact modal, participial and perception expressions in the sea passage remain locally supplied. Cook’s p.143 note18 explicitly requires “give”, not “forgive”, for the final participle.
% mdparagraph:P490:end

\endgroup
\end{document}
